% Manuscript version v16: immutable reproducibility artifact linked.
\documentclass[11pt]{article}

\usepackage[T1]{fontenc}
\usepackage{lmodern}
\usepackage[margin=1in]{geometry}
\usepackage{amsmath,amssymb,mathtools,amsthm}
\usepackage{microtype}
\usepackage[hidelinks]{hyperref}
\usepackage{url}
\usepackage{booktabs}

\allowdisplaybreaks[2]

\newtheorem{theorem}{Theorem}[section]
\newtheorem{proposition}[theorem]{Proposition}
\newtheorem{lemma}[theorem]{Lemma}
\newtheorem{corollary}[theorem]{Corollary}
\newtheorem{remark}[theorem]{Remark}
\DeclareMathOperator{\tr}{tr}
\DeclareMathOperator{\rank}{rank}
\newcommand{\Frob}{\mathrm F}
\newcommand{\R}{\mathbb R}
\newcommand{\C}{\mathbb C}
\newcommand{\cS}{\mathcal S}
\newcommand{\cP}{\mathcal P}
\newcommand{\cD}{\mathcal D}
\newcommand{\cE}{\mathcal E}

\title{A position-weighted refinement for simple zeros of the Riemann zeta function}
\author{Michael Hurtado}
\date{August 12, 2026}

\begin{document}
\maketitle

\begin{abstract}
We prove the unconditional lower bound
\[
 \liminf_{T\to\infty}\frac{N_0^s(T,2T)}{N(T,2T)}
 \ge \frac{655000H_{\rm MT}-1305}{652504}
 =0.6730732086087052768351\ldots,
\]
where
\[
 H_{\rm MT}=\frac32-\frac1{\sqrt2}\cot\frac1{\sqrt2}
 =0.6725007036794116\ldots .
\]
The new ingredient is a position-weighted refinement of the seven-point
Gram-stability argument appearing in recent reproducible work. Earlier
certificates use a uniform pressure on the six consecutive gaps. We observe
that the global double-counting argument depends only on the total pressure,
whereas the local extremal problem depends on its distribution. We exploit
this freedom using the certified vector
\[
 \frac1{10^7}(2714,3733,3553,3553,3733,2714),
 \qquad \sum_{j=1}^6p_j=\frac1{500},
\]
and an Arb/FLINT interval certificate proving the corresponding seven-point
lower bound \(39/10000\). Combined with block length \(m=262\), this yields
the constant above.

For logical independence we also reconstruct the analytic/Gabor input package
from classical ingredients rather than invoking the recent Theorem D as a
black box. The external analytic trust base consists of the Guinand--Weil
explicit formula, Stirling's formula, the Prime Number Theorem, the
Riemann--von Mangoldt formula and standard local zero counting, standard
prime-power estimates, the weighted Montgomery--Vaughan Hilbert inequality,
and Poisson summation. The finite certificate is verified by directed interval
arithmetic.
\end{abstract}

\section{Introduction and statement of novelty}
\label{sec:introduction}

Let \(N(T,2T)\) denote the number of nontrivial zeros of \(\zeta(s)\) with
ordinary ordinate in \((T,2T]\), counted with multiplicity, and let
\(N_0^s(T,2T)\) denote the number of simple zeros on the critical line in the
same interval. The quantitative study of \(N_0^s/N\) goes back to
Montgomery's pair-correlation method and the Montgomery--Taylor optimization.
Recent work has made the relevant finite-compression and zero-side inertia
mechanism unconditional and has produced reproducible improvements beyond the
Montgomery--Taylor constant; see Claude~\cite{Claude26} for the immediate
comparison point and
Baluyot--Goldston--Suriajaya--Turnage-Butterbaugh~\cite{BGSTB24,BGSTB25} and
Goldston--Suriajaya~\cite{GS25,GS26} for surrounding pair-correlation work.
For the optimized Montgomery--Taylor test family we write
\begin{equation}\label{eq:HMT}
 H_{\rm MT}:=\frac32-\frac1{\sqrt2}\cot\frac1{\sqrt2}
 =0.6725007036794116457\ldots .
\end{equation}

Two very recent public computational manuscripts are especially close to the
present argument. Jo's reproducible draft~\cite{Jo26} obtains
\[
 0.6730085279277797613\ldots
\]
using a stability-enhanced rank--trace inequality, a seven-point local
functional, and a uniform pressure on its six gaps. A subsequent repository
maintained by Rademacher~\cite{Rademacher26}, explicitly based on Jo's public
artifact, strengthens the uniform seven-point certificate and obtains
\[
 0.6730213619501665335\ldots .
\]
Jo's repository identifies its research draft as generated by GPT-5.6 Sol.
Rademacher's repository records Jo's artifact as its upstream provenance and
presents the subsequent refinement as AI-generated.  We cite the
human-maintained artifacts and treat model use as provenance rather than
mathematical authorship.

The contribution of the present paper is narrower and explicit. We do not
claim the stability-enhanced rank--trace inequality, the seven-point
architecture, or the uniform-pressure local-to-global argument as new.
Instead we optimize a degree of freedom left fixed in those works. If
\(g_1,\ldots,g_6\) are six consecutive normalized gaps, the global
double-counting cost depends only on
\[
 \beta:=\sum_{j=1}^6p_j,
\]
whereas the local extremal problem depends on the individual \(p_j\).
Consequently the uniform choice can be replaced by a nonuniform
position-dependent pressure without increasing the global cost. The
certified choice used here is
\begin{equation}\label{eq:p}
 p=\frac1{10^7}(2714,3733,3553,3553,3733,2714),
 \qquad \sum_{j=1}^6p_j=\frac1{500}.
\end{equation}
For this vector we certify
\begin{equation}\label{eq:local-cert-intro}
 \mathcal F_p(g_1,\ldots,g_6)\ge\frac{39}{10000}
 \qquad(g_j\ge0).
\end{equation}
After global redistribution and shifted-block averaging, the block length
\(m=262\) yields the theorem below.

\begin{theorem}[Main theorem]\label{thm:main}
With
\[
 H_{\rm MT}=\frac32-\frac1{\sqrt2}\cot\frac1{\sqrt2},
\]
one has unconditionally
\[
 \liminf_{T\to\infty}\frac{N_0^s(T,2T)}{N(T,2T)}
 \ge\frac{655000H_{\rm MT}-1305}{652504}
 =0.6730732086087052768351\ldots .
\]
\end{theorem}

A second contribution is expository and logical rather than a separate
priority claim: we reconstruct the analytic/Gabor package needed by the proof
from classical theorems and explicitly identify the common finite compression
on its zero and prime sides. This makes the proof independent of using the
recent Theorem D as an imported black box.

\paragraph{Priority convention.}
The priority statements above concern the public artifacts inspected while
preparing this version. Jo's repository is the public upstream artifact for
the stability/seven-point draft used here; Rademacher's repository records
that provenance and strengthens its finite certificate. The present novelty
claim is restricted to the nonuniform position-weighted refinement, its
certified pressure vector, and the resulting constant.

\section{The common finite compression and its two representations}
\label{sec:common-compression}

Put
\[
 \ell=\log\frac{T}{2\pi},\quad L=\ell,\quad
 h=\frac{2\pi}{L},\quad
 d=\left\lfloor\frac{T}{h}\right\rfloor,\quad
 \tau_j=T+jh\quad(0\le j<d).
\]
Let $\phi=\phi_L\in C_c^2(\R)$ be the real, even Montgomery--Taylor window
supported on $[-L/2,L/2]$, and define
\[
 \widehat\phi(z):=\int_{\R}\phi(u)e^{-izu}\,du,\qquad
 \mathcal F_+f(z):=\int_{\R}f(u)e^{izu}\,du.
\]
Set
\begin{equation}\label{eq:test-family}
 f_j(u):=\phi(u)e^{-i\tau_j u}.
\end{equation}
Since $\phi$ is even, $\widehat\phi$ is even entire, hence
\begin{equation}\label{eq:test-transform}
 F_j(z):=\mathcal F_+f_j(z)
 =\widehat\phi(\tau_j-z)
 =\widehat\phi(z-\tau_j).
\end{equation}

For a zero $\rho=\beta+i\gamma$ put
\begin{equation}\label{eq:zero-coordinate}
 z_\rho:=-i\left(\rho-\frac12\right)
 =\gamma-i\left(\beta-\frac12\right).
\end{equation}
For the horizontal partner $\rho^\sharp:=1-\overline\rho$,
\begin{equation}\label{eq:partner-coordinate}
 z_{\rho^\sharp}=\overline{z_\rho}.
\end{equation}

Let $W$ denote the Hermitian Weil form in the normalization of
Appendix~\ref{app:inputIII}.  Define the single common finite compression
\begin{equation}\label{eq:G-common}
 G_{jk}:=W(f_j,f_k),\qquad 0\le j,k<d.
\end{equation}
For $c=(c_0,\ldots,c_{d-1})^t$, put
\[
 f_c:=\sum_{j=0}^{d-1}\overline{c_j}\,f_j,\qquad
 \Lambda_\rho(c):=\mathcal F_+f_c(z_\rho)
 =\sum_{j=0}^{d-1}\overline{c_j}\,F_j(z_\rho).
\]
With this convention, which is chosen only to match the standard matrix
quadratic form, one has $c^*Gc=W(f_c,f_c)$ exactly.

\begin{proposition}[Two representations of the common compression]
\label{prop:two-representations}
Let $X=e^L=T/(2\pi)$ and let $\nu_X=\mu+P_X+\Pi_X$ be the density derived
from the classical Weil explicit formula in
Proposition~\ref{aiii:lem:spectral}. Then:

\smallskip
\noindent{\rm (a) Spectral representation.}
\begin{equation}\label{eq:G-spectral-common}
 G_{jk}
 =
 \int_{\R}
 \widehat\phi(\tau-\tau_j)
 \overline{\widehat\phi(\tau-\tau_k)}
 \nu_X(\tau)\,d\tau.
\end{equation}

\smallskip
\noindent{\rm (b) Zero-side representation.}
For every $c\in\C^d$,
\begin{equation}\label{eq:G-zero-common}
 c^*Gc
 =
 \sum_\rho m_\rho\,
 \Lambda_\rho(c)\,
 \overline{\Lambda_{\rho^\sharp}(c)}.
\end{equation}
Grouping under $\rho\leftrightarrow\rho^\sharp$ gives
\begin{align}
 c^*Gc
 &=
 \sum_{\rho=\rho^\sharp}
 m_\rho|\Lambda_\rho(c)|^2 \nonumber\\
 &\quad+
 \sum_{\{\rho,\rho^\sharp\},\,\rho\ne\rho^\sharp}
 2m_\rho\Re\!\left(
 \Lambda_\rho(c)\overline{\Lambda_{\rho^\sharp}(c)}
 \right).
 \label{eq:G-zero-blocks}
\end{align}

\smallskip
\noindent{\rm (c) Interior/exterior decomposition.}
Let $I'=(T-T^{1/2},2T+T^{1/2}]$. Define $A$ by restricting the zero-side
sum to zeros with ordinary ordinates in $I'$, and $E$ by the complementary
sum. Since $\rho\mapsto\rho^\sharp$ preserves ordinary ordinate,
\begin{equation}\label{eq:G-AE-common}
 G=A+E
\end{equation}
entrywise.

\smallskip
\noindent{\rm (d) Normalized evaluation vectors.}
Put
\[
 a=\frac1L\int_{\R}\phi(u)^2\,du,\qquad
 \widehat A=\frac{A}{aL^2}.
\]
For a simple critical-line zero $\rho$ in $I'$, define
\[
 v_\rho:=
 \frac1{\sqrt a\,L}
 \bigl(F_0(\gamma),\ldots,F_{d-1}(\gamma)\bigr)^t .
\]
Then $\Lambda_\rho(c)/(\sqrt a\,L)=c^*v_\rho$, and
\begin{equation}\label{eq:vnorm-common}
 \|v_\rho\|^2
 =
 \frac1{aL^2}\sum_{j=0}^{d-1}
 |\widehat\phi(\gamma-\tau_j)|^2
 \le1.
\end{equation}
Thus
\[
 P_1:=\sum_{\rho\in\cS_1}v_\rho v_\rho^*
\]
is exactly the normalized simple-critical-line part of $\widehat A$, and
$Q':=\widehat A-P_1$ is the pull-back of the remaining zero-side blocks.
\end{proposition}

\begin{proof}
Proposition~\ref{aiii:lem:spectral} and
\eqref{eq:test-transform} give \eqref{eq:G-spectral-common}. The zero-side
form of the same Weil identity is
\[
 W(f,g)=\sum_\rho m_\rho\,
 \mathcal F_+f(z_\rho)\,
 \overline{\mathcal F_+g(\overline{z_\rho})}.
\]
By the coefficient convention above,
$c^*Gc=W(f_c,f_c)$.  Using \eqref{eq:partner-coordinate} with
$f=g=f_c$ therefore proves \eqref{eq:G-zero-common}; grouping paired zeros
proves
\eqref{eq:G-zero-blocks}. Splitting this same zero sum by whether the ordinary
ordinate lies in $I'$ proves \eqref{eq:G-AE-common}.

For a critical-line zero, $z_\rho=\gamma\in\R$. Extending the finite sum in
\eqref{eq:vnorm-common} to the full translated critical lattice and applying
Lemma~\ref{aiv:lem:poisson} at equal arguments (translation of the lattice
changes only nonzero dual phases, whose alias terms vanish by exact support)
gives
\[
 \sum_{j\in\mathbb Z}|\widehat\phi(\gamma-\tau_j)|^2=aL^2.
\]
Hence the finite sum is at most $aL^2$. Finally,
\eqref{eq:G-zero-blocks} identifies the simple part with
$\sum v_\rho v_\rho^*$ and the remainder with the pull-back of the multiple
and off-line blocks.
\end{proof}

\begin{remark}[Dependency graph]
Input I uses the zero-side representation above, Input II uses $G=A+E$,
Input III uses the spectral representation, and Input IV supplies the
critical-lattice identity used in \eqref{eq:vnorm-common}.  Thus all four
modules refer to one test family and one matrix.
\end{remark}


\section{Analytic/Gabor package and the stability refinement}\label{sec:analytic}

Put
\[
 \ell:=\log\frac{T}{2\pi},\qquad I=(T,2T],\qquad
 D_0=T^{1/2},\qquad I'=(T-D_0,2T+D_0],
\]
and set $L=\ell$.  Let $\phi$ be the real, even, compactly tapered
Montgomery--Taylor window used below, and define
\[
 a:=\frac1L\int_{\R}\phi(u)^2\,du,\qquad
 \widehat G:=\frac{G}{aL^2},\qquad
 \widehat A:=\frac{A}{aL^2}.
\]

The analytic work needed by the stability argument is summarized in the next
proposition.  Its four assertions are proved independently in
Appendices~\ref{app:inputI}--\ref{app:inputIV}.

\begin{proposition}[Analytic/Gabor package]\label{prop:analytic-package}
Let the zeros in $I'$ be partitioned into simple critical-line zeros
$\cS_1$, distinct multiple critical-line zeros $\cS_2$, and unordered
horizontal pairs
\[
 \{\rho,\rho^\sharp\},\qquad \rho^\sharp=1-\overline\rho,
\]
off the critical line.  Write $s_1=\#\cS_1$, $s_2=\#\cS_2$ and
$p=\#\cP$.  Then the following hold.

\smallskip
\noindent{\rm (I) Zero-side inertia.}
\begin{equation}\label{eq:Nlower}
 N(I')\ge s_1+2s_2+2p.
\end{equation}
For the normalized sampled vectors $v_\rho$ associated with
$\rho\in\cS_1$, let $V$ have these vectors as columns and put
\[
 P_1=VV^*,\qquad M=V^*V,\qquad Q'=\widehat A-P_1.
\]
Then
\begin{equation}\label{eq:index}
 n_+(Q')\le s_2+p.
\end{equation}

\smallskip
\noindent{\rm (II) Exterior-zero tail.}
Writing $\widehat A=\widehat G-\widehat E$,
\begin{equation}\label{eq:tail-input}
 4\tr\widehat A-\|\widehat A\|_{\Frob}^2
 \ge
 4\tr\widehat G-\|\widehat G\|_{\Frob}^2-o(N(T,2T)).
\end{equation}
In fact the proof gives $\|\widehat E\|_1\ll T^{-1/2}$.

\smallskip
\noindent{\rm (III) Optimized first and second moments.}
\begin{equation}\label{eq:TDmoments}
 \tr\widehat G=N(T,2T)(1+o(1)),\qquad
 \|\widehat G\|_{\Frob}^2=
 \left(\frac1{c_1^*}+o(1)\right)N(T,2T),
\end{equation}
where
\begin{equation}\label{eq:c1}
 \frac1{c_1^*}
 =
 \frac12+\frac1{\sqrt2}\cot\frac1{\sqrt2},
 \qquad
 2-\frac1{c_1^*}=H_{\rm MT}.
\end{equation}

\smallskip
\noindent{\rm (IV) Critical-lattice Poisson--Gabor limit.}
At mesh $h=2\pi/L$, Poisson summation gives the infinite-grid identity
without nonzero aliases.  The tapered Fourier decay controls finite-grid
truncation, and the normalized overlap kernel converges to the compact
Montgomery--Taylor kernel used in Section~\ref{sec:kernel}.  The integrated
second-moment end effect is
\[
 O\!\left(L\ell\log L\,(\ell^2+X)\right),
 \qquad X=e^L\asymp T,
\]
which is $o(TL\ell^2)$ at $L=\ell$.
\end{proposition}

\begin{proof}
Assertion (I) is Theorem~\ref{ai:thm:InputI}; assertion (II) is
Theorem~\ref{aii:thm:II}; assertion (III) is proved in
Appendix~\ref{app:inputIII}, culminating in
\eqref{aiii:eq:cstar}; and assertion (IV) is
Theorem~\ref{aiv:thm:IV}.  These appendices use only the classical trust base
listed in Section~\ref{sec:trust}.
\end{proof}

Define the convex function
\begin{equation}\label{eq:Psi}
 \Psi(t)=
 \begin{cases}
  (t-1)^2,&0\le t\le2,\\
  2t-3,&t\ge2,
 \end{cases}
 \qquad
 \cD(M):=\tr\Psi(M)
\end{equation}
for $M\succeq0$.

\begin{lemma}[Stability-enhanced rank--trace inequality]\label{lem:stable}
Let $V$ be a $d\times r$ matrix whose columns have norm at most one.  Put
$P=VV^*\succeq0$ and $M=V^*V$.  If $Q$ is Hermitian and $n_+(Q)\le b$, then
\begin{equation}\label{eq:stable}
 \|P+Q\|_{\Frob}^2
 \ge 4\tr(P+Q)-3r-4b+\cD(M).
\end{equation}
\end{lemma}

\begin{proof}
Write $Q=Q_+-Q_-$ for its positive and negative parts.  They are positive
semidefinite, have orthogonal supports, and $\operatorname{rank}Q_+\le b$.
Since $\tr((P-Q_-)Q_+)=\tr(PQ_+)\ge0$,
\[
 \|P+Q\|_{\Frob}^2
 \ge \|P-Q_-\|_{\Frob}^2+\|Q_+\|_{\Frob}^2.
\]
If $q>0$, then $q^2\ge4q-4$; summing over the at most $b$ positive
eigenvalues of $Q$ gives
\begin{equation}\label{eq:qplus}
 \|Q_+\|_{\Frob}^2\ge4\tr Q_+-4b.
\end{equation}

Let $p_1\ge\cdots\ge p_r\ge0$ be the eigenvalues of $M=V^*V$
(counted with multiplicity, hence including any zeros).  The nonzero
eigenvalues of $P=VV^*$ are the same; extend the eigenvalue list of $P$ by
zeros as necessary.  Likewise write
$n_1\ge n_2\ge\cdots\ge0$ for the eigenvalues of $Q_-$ and extend either
list by zeros to a common length.  Von Neumann's trace inequality then implies
\[
 \tr(PQ_-)\le\sum_{i\le r}p_i n_i.
\]
Consequently
\[
 \|P-Q_-\|_{\Frob}^2+4\tr Q_-
 \ge\sum_{i\le r}\bigl((p_i-n_i)^2+4n_i\bigr).
\]
For fixed $p\ge0$,
\[
 \min_{n\ge0}\bigl((p-n)^2+4n\bigr)
 =
 \begin{cases}
 p^2,&0\le p\le2,\\
 4p-4,&p\ge2,
 \end{cases}
 =2p-1+\Psi(p).
\]
Therefore
\begin{equation}\label{eq:qminus}
 \|P-Q_-\|_{\Frob}^2
 \ge2\tr P-r+\cD(M)-4\tr Q_-.
\end{equation}
Combining \eqref{eq:qplus} and \eqref{eq:qminus},
\[
 \|P+Q\|_{\Frob}^2
 \ge4\tr(P+Q)-2\tr P-r-4b+\cD(M).
\]
Finally $\tr P$ is the sum of the squared column norms of $V$, so
$\tr P\le r$, proving \eqref{eq:stable}.
\end{proof}

Before stating the global consequence, define the retained central set
\begin{equation}\label{eq:Scentral}
 \cS^\circ:=\{\rho\in\cS_1:T<\gamma_\rho\le2T\text{ and }\rho
 \text{ lies outside the two boundary strips of Section~\ref{sec:kernel}}\}.
\end{equation}
Let $S^\circ:=\#\cS^\circ$ and let
\begin{equation}\label{eq:Mcircdef}
 M^\circ:=\bigl(\langle v_\rho,v_{\rho'}\rangle\bigr)_{\rho,\rho'\in\cS^\circ}.
\end{equation}
Thus $M^\circ$ is a principal submatrix of $M$ and contains only simple
critical-line zeros whose ordinates lie in $(T,2T]$ and for which the compact
kernel approximation is uniform.

\begin{corollary}[Global Gram-defect inequality]\label{cor:global}
Then
\begin{equation}\label{eq:global}
 N_0^s(T,2T)
 \ge H_{\rm MT}N(T,2T)+\cD(M^\circ)-o(N(T,2T)).
\end{equation}
\end{corollary}

\begin{proof}
Apply Lemma~\ref{lem:stable} to $\widehat A=P_1+Q'$ with $r=s_1$ and
$b=s_2+p$.  By \eqref{eq:Nlower} and \eqref{eq:index},
\begin{align*}
 \|\widehat A\|_{\Frob}^2
 &\ge4\tr\widehat A-3s_1-4s_2-4p+\cD(M)\\
 &\ge4\tr\widehat A-s_1-2N(I')+\cD(M).
\end{align*}
Hence
\begin{equation}\label{eq:s1}
 s_1\ge4\tr\widehat A-\|\widehat A\|_{\Frob}^2-2N(I')+\cD(M).
\end{equation}
Apply the independently proved tail estimate~\eqref{eq:tail-input}.  Pinching $M$ to the
retained central principal block and its complement does not increase
$\tr\Psi$: pinching is a random-unitary average, and the spectral trace
functional $X\mapsto\tr\Psi(X)$ is convex and unitarily invariant.  No
operator convexity of $\Psi$ is required.  Thus
$\cD(M)\ge\cD(M^\circ)$.  Moreover,
\[
 N(I')=N(T,2T)+o(N(T,2T)).
\]
Because $s_1$ counts simple critical-line zeros in the enlarged interval
$I'$, whereas $N_0^s(T,2T)$ counts only those in $I$, we also use
\begin{equation}\label{eq:s1-central}
 s_1\le N_0^s(T,2T)+N(I'\setminus I).
\end{equation}
The standard local zero-count estimate gives
\[
 N(I'\setminus I)\ll D_0\log T=T^{1/2}\log T=o(N(T,2T)),
\]
and hence
\begin{equation}\label{eq:N0-from-s1}
 N_0^s(T,2T)\ge s_1-o(N(T,2T)).
\end{equation}
Substitution of \eqref{eq:tail-input}, \eqref{eq:TDmoments}, and
\eqref{eq:c1} into \eqref{eq:s1}, followed by \eqref{eq:N0-from-s1}, gives
\eqref{eq:global}.
\end{proof}

\section{The limiting overlap kernel}\label{sec:kernel}

Put
\[
 h:=\frac{2\pi}{L},\qquad
 x_\rho:=\frac{\gamma_\rho-T}{h}
 =\frac{L(\gamma_\rho-T)}{2\pi}.
\]
Define
\begin{align}
 K(x)
 &:=\int_{-1/2}^{1/2}\cos(\sqrt2\,t)\cos(2\pi xt)\,dt \label{eq:Kint}\\
 &=\frac{\sin(\pi x-1/\sqrt2)}{2\pi x-\sqrt2}
  +\frac{\sin(\pi x+1/\sqrt2)}{2\pi x+\sqrt2}, \label{eq:K}\\
 k(x)&:=\frac{K(x)}{K(0)},\qquad
 w(x):=k(x)^2,\qquad
 K(0)=\sqrt2\sin(1/\sqrt2). \label{eq:kw}
\end{align}
The apparent singularities in \eqref{eq:K} are removable.

\begin{lemma}[Compact overlap limit]\label{lem:kernel}
Fix $R_0>0$.  Delete simple zeros lying in boundary strips of normalized
width $L^2$ at the two ends of the sampling grid.  The number deleted is
$o(N(T,2T))$.  Uniformly for retained simple zeros satisfying
$|x_\rho-x_{\rho'}|\le R_0$,
\begin{equation}\label{eq:kernel-limit}
 \langle v_\rho,v_{\rho'}\rangle
 =k(x_\rho-x_{\rho'})+o(1).
\end{equation}
\end{lemma}

\begin{proof}
By Appendix~\ref{app:inputIV}, the independently proved Poisson--Gabor
sampling identity, applied to the full grid, gives
\[
 \frac1{aL^2}\sum_{j\in\mathbb Z}
 \widehat\phi(\gamma-\tau_j)\widehat\phi(\gamma'-\tau_j)
 =\frac{\Phi(\gamma-\gamma')}{aL},
 \qquad \Phi=\widehat{\phi^2}.
\]
We record explicitly the pointwise truncation estimate needed here.  The
Fourier decay of the fixed-width tapered window gives, uniformly in $L$,
\[
 |\widehat\phi(r)|\ll
 \min\left\{L,\frac1{|r|},\frac1{r^2}\right\}.
\]
Suppose that both points lie at normalized distance at least $L^2$ from the
finite grid ends and that $|x_\rho-x_{\rho'}|\le R_0$.  If an omitted grid
index lies $n\ge L^2$ grid steps beyond the nearer end, then
\[
 |\gamma-\tau_j|\asymp \frac nL,\qquad
 |\gamma'-\tau_j|\asymp \frac nL
\]
uniformly for fixed $R_0$.  Hence
\[
 |\widehat\phi(\gamma-\tau_j)
   \widehat\phi(\gamma'-\tau_j)|
 \ll \frac{L^4}{n^4}.
\]
Summing the two omitted tails gives
\[
 \sum_{n\ge L^2}\frac{L^4}{n^4}\ll L^{-2}.
\]
For the optimized window, dominated convergence also gives
\[
 a=\frac1L\int_{\R}\phi(u)^2\,du
   =K(0)+o(1),
 \qquad K(0)=\sqrt2\sin(1/\sqrt2)>0,
\]
so in particular $a\asymp1$.  After multiplication by the Gram normalization
$1/(aL^2)$, the omitted tails therefore contribute $O(L^{-4})=o(1)$,
uniformly for bounded normalized separation.  Thus the
full-grid Poisson--Gabor identity may be replaced by the finite sampled inner
product with a uniform $o(1)$ error.

For $\gamma-\gamma'=hx$ and fixed $x$, substitute $u=Lt$:
\[
 \frac{\Phi(hx)}{aL}
 =
 \frac{\displaystyle\int_{-1/2}^{1/2}
       \phi(Lt)^2e^{-2\pi ixt}\,dt}
      {\displaystyle\int_{-1/2}^{1/2}\phi(Lt)^2\,dt}.
\]
The squared tapered window converges in $L^1$ to
$\cos(\sqrt2\,t)\mathbf1_{[-1/2,1/2]}(t)$; the rescaled taper occupies
length $O(L^{-1})$.  Hence the last quotient converges to $K(x)/K(0)$,
uniformly for $x$ in compact sets.  Combining this with the finite-grid tail
bound proves \eqref{eq:kernel-limit}.

Each deleted boundary strip has ordinate length $O(L)$, and the standard local
zero-count estimate $N(t+1)-N(t)\ll\log(t+3)$ shows that only $O(L^2)=o(N)$
zeros are removed.
\end{proof}

\begin{lemma}[Block defect]\label{lem:block}
If $G\succeq0$ is Hermitian, then
\begin{equation}\label{eq:block}
 \tr\Psi(G)\ge
 \min\left\{1,\;2\sum_{i<j}|G_{ij}|^2\right\}.
\end{equation}
\end{lemma}

\begin{proof}
If all eigenvalues of $G$ lie in $[0,2]$, then
$\Psi(G)=(G-I)^2$ and therefore
\[
 \tr\Psi(G)=\|G-I\|_{\Frob}^2
 =\sum_i(G_{ii}-1)^2+2\sum_{i<j}|G_{ij}|^2
 \ge 2\sum_{i<j}|G_{ij}|^2.
\]
If some eigenvalue $\lambda>2$, then its contribution is
$\Psi(\lambda)=2\lambda-3>1$, while all other contributions are nonnegative.
This proves \eqref{eq:block}.
\end{proof}

\section{Position-weighted seven-point stability}\label{sec:local}

Let $g_1,\ldots,g_6\ge0$ be six consecutive gaps.  For
$1\le r\le6$ and $1\le i\le7-r$, write
\[
 s_{i,r}=g_i+\cdots+g_{i+r-1}
\]
and define
\begin{equation}\label{eq:E7}
 \cE_7(g):=
 \sum_{r=1}^6\frac{2}{7-r}
 \sum_{i=1}^{7-r}w(s_{i,r}).
\end{equation}
For $p_j\ge0$ put
\begin{equation}\label{eq:Fp}
 \mathcal F_p(g):=\sum_{j=1}^6p_jg_j+\cE_7(g).
\end{equation}

\begin{proposition}[Pressure redistribution]\label{prop:redist}
Let $\beta=\sum_{j=1}^6p_j$.  If
\[
 \mathcal F_p(g)\ge\delta
 \qquad(g\in\R_{\ge0}^6),
\]
then every ordered configuration $y_1<\cdots<y_m$, $m\ge7$, satisfies
\begin{equation}\label{eq:block-energy-general}
 E_m+\beta(y_m-y_1)\ge\delta(m-6),
 \qquad
 E_m:=2\sum_{1\le i<j\le m}w(y_j-y_i).
\end{equation}
\end{proposition}

\begin{proof}
Sum the local inequality over the $m-6$ consecutive seven-point windows.
A fixed pair separated by $r$ gaps can appear in at most $7-r$ windows and
has coefficient $2/(7-r)$ in each appearance; hence its total coefficient is
at most $2$.  A fixed global gap can occupy each of the six relative gap
positions at most once, so its accumulated pressure is at most
$\sum_jp_j=\beta$.  Therefore
\[
 (m-6)\delta
 \le\sum_{\text{windows}}\mathcal F_p
 \le E_m+\beta\sum_{j=1}^{m-1}(y_{j+1}-y_j),
\]
which is \eqref{eq:block-energy-general}.
\end{proof}

\begin{proposition}[Computer-assisted seven-point inequality]\label{prop:cert}
For the pressure vector \eqref{eq:p}, every $g_1,\ldots,g_6\ge0$ satisfies
\begin{align}
 &\frac{2714(g_1+g_6)+3733(g_2+g_5)+3553(g_3+g_4)}{10^7}\nonumber\\
 &\qquad+
 \sum_{r=1}^6\frac{2}{7-r}
 \sum_{i=1}^{7-r}w(g_i+\cdots+g_{i+r-1})
 \ge\frac{39}{10000}.
 \label{eq:cert}
\end{align}
\end{proposition}

\begin{proof}[Computer-assisted proof]
The archived verifier performs exhaustive interval branch-and-bound with
Arb/FLINT through \texttt{python-flint}.  The pressure coefficients, target,
and grid endpoints are represented by exact rationals.  The verifier
reconstructs rigorous Arb enclosures for $w$ and its required derivatives from
\eqref{eq:K}--\eqref{eq:kw}.  The grid mesh is $1/4000$.

Boxes are discarded only after one of three rigorous lower bounds proves the
target: direct interval evaluation, the positive pressure bound, or a tangent
bound after positive definiteness of the interval Hessian has been certified.
Otherwise the widest coordinate is bisected.  Reaching an uncertified terminal
grid cell raises an exception; thus a successful run closes the entire search
tree.

The frozen 256-bit run reports \texttt{verified=true}, with $1{,}119{,}372$
visited nodes, $559{,}038$ splits, and maximum depth $39$.  The same frozen
verifier also succeeds at 128-bit working precision.  The complete 256-bit
artifact is publicly archived in the immutable release
\cite{Hurtado26Artifact}.  Appendix~\ref{app:cert} records the hashes and
environment.
\end{proof}

\begin{corollary}\label{cor:energy}
For every $m\ge7$ and $y_1<\cdots<y_m$,
\begin{equation}\label{eq:energy}
 E_m+\frac1{500}(y_m-y_1)
 \ge\frac{39}{10000}(m-6).
\end{equation}
\end{corollary}

\begin{proof}
Apply Proposition~\ref{prop:redist} to Proposition~\ref{prop:cert} and use
$\sum_jp_j=1/500$ from \eqref{eq:p}.
\end{proof}

\section{Block stability}\label{sec:blockstability}

Take
\begin{equation}\label{eq:mA}
 m=262,\qquad
 A_0:=\frac{39}{10000}(m-6)=\frac{624}{625}<1.
\end{equation}
This is the largest integer $m$ for which $(39/10000)(m-6)<1$.

Let $B$ be a block of $m$ consecutive retained simple zeros, let
$\operatorname{span}(B)$ denote the difference between its extreme normalized
ordinates, and let $G_B$ be its Gram matrix.

\begin{lemma}[Uniform block stability]\label{lem:blockstab}
Uniformly over such blocks,
\begin{equation}\label{eq:blockstab}
 \cD(G_B)+\frac1{500}\operatorname{span}(B)\ge A_0-o(1).
\end{equation}
\end{lemma}

\begin{proof}
If $\operatorname{span}(B)/500\ge A_0$, the claim is immediate because
$\cD(G_B)\ge0$.  Otherwise
$\operatorname{span}(B)<500A_0<500$.  Since $m$ is fixed, all
$\binom{m}{2}$ pair separations lie in one fixed compact interval.  By
Lemma~\ref{lem:kernel},
\[
 2\sum_{i<j}|(G_B)_{ij}|^2=E_m+o(1)
\]
uniformly over the blocks in this branch.  Lemma~\ref{lem:block} gives
\[
 \cD(G_B)\ge\min\{1,E_m+o(1)\}.
\]
Corollary~\ref{cor:energy} gives
\[
 E_m\ge A_0-\frac1{500}\operatorname{span}(B).
\]
The right-hand side lies in $(0,A_0]$ and $A_0<1$.  Hence
\[
 \cD(G_B)
 \ge A_0-\frac1{500}\operatorname{span}(B)-o(1),
\]
which proves \eqref{eq:blockstab}.  Uniformity of the error follows from the
fixed block size and the compact separation range.
\end{proof}

\section{Shifted block pinching}\label{sec:pinching}

Let
\[
 x_1<\cdots<x_{S^\circ}
\]
be the normalized ordinates of the set $\cS^\circ$ defined in
\eqref{eq:Scentral}.  By the boundary estimate in Lemma~\ref{lem:kernel},
\begin{equation}\label{eq:S}
 S^\circ=N_0^s(T,2T)-o(N(T,2T)).
\end{equation}

For each of the $m$ offsets modulo $m$, partition the ordered points into full
consecutive blocks of size $m$, leaving only $O(m)$ endpoint points.  Pinching
$M^\circ$ to the full principal blocks and the remainder yields
\begin{equation}\label{eq:pinch}
 \cD(M^\circ)\ge\sum_B\cD(G_B).
\end{equation}
Indeed, if $\mathcal P$ denotes this pinching, then for the finite group
$\mathcal U$ of blockwise sign unitaries one has the random-unitary
representation
\[
 \mathcal P(X)=\frac1{|\mathcal U|}
 \sum_{U\in\mathcal U}UXU^*.
\]
Since the spectral trace functional
$F(X):=\tr\Psi(X)$ is convex and unitarily invariant,
\[
 F(\mathcal P(X))
 \le\frac1{|\mathcal U|}\sum_{U\in\mathcal U}F(UXU^*)
 =F(X).
\]
The pinched matrix is the direct sum of the full block Gram matrices and a
remainder block.  Since $\Psi\ge0$, the latter has nonnegative defect and may
be discarded, proving \eqref{eq:pinch}.  No operator convexity of $\Psi$ is
used.

For a fixed offset, sum \eqref{eq:blockstab} over its full blocks; then average
over all $m$ offsets.  The average number of full blocks is
$S^\circ/m+O(1)$.  Each consecutive gap $x_{j+1}-x_j$ belongs to a block span
for at most $m-1$ offsets.  Moreover,
\begin{equation}\label{eq:span}
 x_{S^\circ}-x_1
 \le\frac{T}{h}
 =\frac{LT}{2\pi}
 =N(T,2T)+o(N(T,2T))
\end{equation}
by the Riemann--von Mangoldt formula.  Since the error in
Lemma~\ref{lem:blockstab} is uniform and $m$ is fixed,
\begin{equation}\label{eq:defectglobal}
 \cD(M^\circ)
 \ge
 \frac{A_0}{m}N_0^s(T,2T)
 -\frac{m-1}{500m}N(T,2T)
 -o(N(T,2T)).
\end{equation}
For \eqref{eq:mA},
\begin{equation}\label{eq:defectnumbers}
 \cD(M^\circ)
 \ge
 \frac{312}{81875}N_0^s(T,2T)
 -\frac{261}{131000}N(T,2T)
 -o(N(T,2T)).
\end{equation}

\begin{proof}[Proof of Theorem~\ref{thm:main}]
Insert \eqref{eq:defectnumbers} into Corollary~\ref{cor:global}:
\[
 N_0^s
 \ge H_{\rm MT}N
 +\frac{312}{81875}N_0^s
 -\frac{261}{131000}N-o(N).
\]
Therefore
\[
 \frac{81563}{81875}N_0^s
 \ge
 \left(H_{\rm MT}-\frac{261}{131000}\right)N-o(N).
\]
Divide by $N=N(T,2T)$ and pass to the lower limit:
\begin{align*}
 \liminf_{T\to\infty}\frac{N_0^s(T,2T)}{N(T,2T)}
 &\ge
 \frac{H_{\rm MT}-261/131000}{81563/81875}\\
 &=\frac{655000H_{\rm MT}-1305}{652504}\\
 &=0.6730732086087052768351\ldots .
\end{align*}
\end{proof}

\section{Comparison and structural interpretation}

Jo's upstream reproducible draft~\cite{Jo26} uses
$p_j=1/3000$ and certifies $19/5000=0.0038$, leading to
$0.6730085279277797\ldots$.  Rademacher's subsequent reproducible
refinement~\cite{Rademacher26} keeps the same uniform pressure and strengthens
the certified local value to $191/50000=0.00382$, leading to
$0.6730213619501665\ldots$.  Our vector has the same total pressure
$1/500$ but certifies $39/10000=0.0039$.

The natural finite-dimensional optimization exposed by
Proposition~\ref{prop:redist} is
\begin{equation}\label{eq:minimax}
 \sup_{\substack{p_j\ge0\\\sum p_j=1/500}}
 \inf_{g\in\R_{\ge0}^6}
 \left(\sum_{j=1}^6p_jg_j+\cE_7(g)\right).
\end{equation}
The uniform vector is one feasible point of this problem.  The strongest reproduced \emph{certificate} for that uniform vector among
the two artifacts just cited is Rademacher's
$191/50000=0.00382$, whereas Proposition~\ref{prop:cert} gives a certified
value $39/10000=0.0039$ for the position-weighted vector \eqref{eq:p}.
This comparison proves that the new \emph{certified lower bound} is stronger;
by itself it does not prove that the exact local infimum for the uniform
vector is below $0.0039$.  Accordingly, we make no claim here that uniform
pressure is globally suboptimal, nor that \eqref{eq:p} solves
\eqref{eq:minimax}.

A larger optimization could also redistribute pair-energy coefficients
$c_{r,i}$ subject to the double-counting budgets
$\sum_{i=1}^{7-r}c_{r,i}\le2$.  We leave that problem open.

\section{Scope and trust base}\label{sec:trust}

\noindent\textbf{Logical status of the theorem.}
The common compression is fixed in Section~\ref{sec:common-compression}, and the analytic/Gabor package used in the main argument is proved in
Appendices~\ref{app:inputI}--\ref{app:inputIV}; it is not imported from
Theorem~D of~\cite{Claude26}.  The analytic external trust base consists of
standard classical results: the Weil explicit formula for $\zeta$ in the
Fourier normalization stated in Appendix~\ref{app:inputIII}, Stirling's
formula, the Prime Number Theorem, the Riemann--von Mangoldt formula and its standard local zero-count consequence, standard Chebyshev and prime-power estimates, the weighted
Montgomery--Vaughan Hilbert inequality, and Poisson summation.  The
Montgomery--Vaughan paper is correctly cited as
\emph{Hilbert's inequality}, J. London Math. Soc. (2) \textbf{8} (1974),
73--82.

The role of~\cite{Claude26} in this version is comparison and normalization
cross-checking only.  No proposition from that preprint is invoked as a
premise of Theorem~\ref{thm:main}.  Likewise, the recent pair-correlation
papers cited in the introduction provide context; the particular
first/second-moment route used here is reconstructed directly in
Appendix~\ref{app:inputIII}.

Proposition~\ref{prop:cert} is computer-assisted.  Its trust base includes the
immutable artifact release~\cite{Hurtado26Artifact}, the frozen verifier,
\texttt{python-flint}, FLINT/Arb and their dependencies, and the correctness
of the execution environment.  The 128- and 256-bit runs are
replications of the same implementation at different working precisions, not
algorithmically independent verifiers.

The pressure vector \eqref{eq:p} is an explicit certified improvement, not an
optimality theorem.  Thus the paper is self-contained relative to the
classical analytic results just listed and to the stated Arb/FLINT
certificate; it does not claim to reprove those classical theorems from first
principles.

\appendix


\section{Independent proof of Input I: zero-side inertia}\label{app:inputI}
This appendix proves the algebraic finite-compression statement used in
Proposition~\ref{prop:analytic-package}.  The horizontal partner
$\rho^\sharp=1-\overline\rho$ is the composition of conjugation with the
functional-equation symmetry; it preserves ordinary ordinate.
\subsection{Scope}

Let $I'$ be a finite ordinate interval. Since
$\rho\mapsto1-\overline\rho$ preserves ordinary ordinate, the zeros with
ordinates in $I'$ form a pairing-invariant finite multiset $Z(I')$ of
nontrivial zeta zeros whose ordinates lie in $I'$.  We partition the
\emph{distinct} zeros into
\[
 \cS_1=\{\text{simple critical-line points}\},\qquad
 \cS_2=\{\text{multiple critical-line points}\},
\]
and a set $\cP$ containing one representative from each off-line pair
\[
 \{\rho,\rho^\sharp\},
 \qquad
 \rho^\sharp:=1-\overline\rho .
\]
Write
\[
 s_1=\#\cS_1,\qquad s_2=\#\cS_2,\qquad p=\#\cP.
\]

The only zeta-specific fact used in the inertia calculation is the combined
functional-equation/conjugation symmetry, which preserves multiplicity:
\[
 m_{\rho^\sharp}=m_\rho.
\]
The finite zero-side Hermitian form obtained by compressing Weil's form to a
finite test family can then be written as a pull-back of the block form
described below.  The present note verifies the resulting inertia statement
from scratch.

\subsection{Positive index under pull-back}

For a Hermitian form $H$ on a finite-dimensional complex vector space, let
$n_+(H)$ denote its positive index: the maximal dimension of a subspace on
which $H$ is positive definite.

\begin{lemma}[Pull-back monotonicity]\label{ai:lem:pullback}
Let $H$ be a Hermitian form on $\C^m$, and let
$A:U\to\C^m$ be either complex-linear or conjugate-linear.  Then
\[
 n_+(H\circ A)\le n_+(H).
\]
\end{lemma}

\begin{proof}
Let $U_0\subset U$ be a complex subspace on which $H\circ A$ is positive
definite.  Then $A|_{U_0}$ is injective.  If $A$ is conjugate-linear,
$A(U_0)$ is still a complex subspace, because
$\alpha A(u)=A(\overline\alpha\,u)$ for every $\alpha\in\C$.
Moreover $H$ is positive definite on $A(U_0)$.  Hence
\[
 \dim_{\C}U_0=\dim_{\C}A(U_0)\le n_+(H).
\]
Taking the maximum over $U_0$ proves the claim.
\end{proof}

\begin{corollary}[Subadditivity]\label{ai:cor:subadd}
For Hermitian forms $H_1,H_2$ on the same space,
\[
 n_+(H_1+H_2)\le n_+(H_1)+n_+(H_2).
\]
\end{corollary}

\begin{proof}
Apply Lemma~\ref{ai:lem:pullback} to $H_1\oplus H_2$ under the diagonal map
$u\mapsto(u,u)$.
\end{proof}

\subsection{The off-line hyperbolic block}

\begin{lemma}[Signature of one horizontal symmetry pair]\label{ai:lem:block}
Let $m>0$.  On $\C^2$ define
\[
 H_m(x,y)=2m\,\operatorname{Re}(x\overline y).
\]
Then
\[
 H_m(x,y)
 =
 \begin{pmatrix}\overline x&\overline y\end{pmatrix}
 \begin{pmatrix}0&m\\m&0\end{pmatrix}
 \binom{x}{y},
\]
and the form has signature $(1,1)$.
\end{lemma}

\begin{proof}
The Hermitian matrix
\[
 B_m=\begin{pmatrix}0&m\\m&0\end{pmatrix}
\]
has characteristic polynomial $\lambda^2-m^2$, hence eigenvalues
$m$ and $-m$.  Equivalently, with
\[
 u=\frac{x+y}{\sqrt2},\qquad v=\frac{x-y}{\sqrt2},
\]
one has
\[
 H_m(x,y)=m|u|^2-m|v|^2.
\]
Thus $n_+(H_m)=n_-(H_m)=1$.
\end{proof}

This is the entire source of the ``one positive direction per off-line
pair'' rule.  No limiting argument and no assumption concerning the distance
of the pair from the critical line is involved.

\subsection{The zero-side block form}

For each distinct critical-line zero $\rho$, let $m_\rho$ be its
multiplicity.  Introduce one evaluation coordinate $z_\rho\in\C$.  For each
off-line representative $\rho\in\cP$, introduce two coordinates
$(z_\rho,z_{\rho^\sharp})$.

Define the Hermitian form
\begin{align}
 \mathcal Q(z)
 &=
 \sum_{\rho\in\cS_1\cup\cS_2}
 m_\rho |z_\rho|^2 \nonumber\\
 &\quad+
 \sum_{\rho\in\cP}
 2m_\rho\operatorname{Re}
 \bigl(z_\rho\overline{z_{\rho^\sharp}}\bigr).
 \label{ai:eq:Q}
\end{align}

\begin{proposition}[Exact block inertia]\label{ai:prop:Qinertia}
The form $\mathcal Q$ is an orthogonal direct sum of
$s_1+s_2$ positive one-dimensional blocks and $p$ hyperbolic two-dimensional
blocks.  Consequently
\[
 n_+(\mathcal Q)=s_1+s_2+p.
\]
If the simple critical-line blocks are removed, the remaining form
$\mathcal Q'$ satisfies
\[
 n_+(\mathcal Q')=s_2+p.
\]
\end{proposition}

\begin{proof}
Every critical-line block is $m_\rho|z|^2$ with $m_\rho>0$, and hence has
positive index one.  Every off-line pair block has signature $(1,1)$ by
Lemma~\ref{ai:lem:block}.  Positive index is additive under orthogonal direct
sum, proving both assertions.
\end{proof}

\subsection{Compression by evaluation}

Let $U=\C^d$ be the coefficient space of the finite family of test functions
used to form the compression.  For each zero parameter appearing on the zero side, let
$\ell_\rho$ denote the conjugate-linear coefficient functional
$c\mapsto\Lambda_\rho(c)$ from
Section~\ref{sec:common-compression}.  Define
\[
 \mathfrak E(c):=\bigl(\ell_\rho(c)\bigr)_\rho .
\]
The map $\mathfrak E$ is conjugate-linear in the displayed coefficient
convention; Lemma~\ref{ai:lem:pullback} was formulated to cover this case
directly.

By Proposition~\ref{prop:two-representations}, the normalized finite
zero-side matrix $\widehat A$ is exactly the matrix of the pull-back form
\[
 c\longmapsto \mathcal Q(\mathfrak E c).
\]
Thus this inertia calculation applies to the same matrix $\widehat A$ used in
the main argument.

For a simple critical-line zero $\rho\in\cS_1$, use the normalized
evaluation vector $v_\rho\in\C^d$ defined in
Proposition~\ref{prop:two-representations}; thus the normalized evaluation is
$c^*v_\rho$.
Proposition~\ref{prop:two-representations} proves for this same family that
\[
 \|v_\rho\|\le1.
\]
Set
\[
 P_1=\sum_{\rho\in\cS_1}v_\rho v_\rho^*,
 \qquad
 M=V^*V,
\]
where $V$ has the $v_\rho$ as columns, and put
\[
 Q'=\widehat A-P_1.
\]

\begin{theorem}[Independent Input I]\label{ai:thm:InputI}
With the preceding notation,
\[
 P_1\succeq0,\qquad
 \rank P_1\le s_1,\qquad
 \tr P_1\le s_1,
\]
and
\[
 n_+(Q')\le s_2+p.
\]
In particular, the finite compression cannot acquire more than one positive
direction from each multiple critical-line point and each off-line
horizontal symmetry pair after the simple critical-line contribution has
been separated.
\end{theorem}

\begin{proof}
Each summand $v_\rho v_\rho^*$ is positive semidefinite and rank at most one,
so
\[
 P_1\succeq0,\qquad \rank P_1\le s_1.
\]
Moreover
\[
 \tr P_1
 =
 \sum_{\rho\in\cS_1}\|v_\rho\|^2
 \le s_1.
\]

After deleting the simple critical-line coordinates from the block form
$\mathcal Q$, the remaining form is $\mathcal Q'$ of
Proposition~\ref{ai:prop:Qinertia}, with
\[
 n_+(\mathcal Q')=s_2+p.
\]
The matrix $Q'$ on coefficient space is the pull-back of this form under the
corresponding restricted evaluation map.  Lemma~\ref{ai:lem:pullback} therefore
gives
\[
 n_+(Q')\le n_+(\mathcal Q')=s_2+p.
\]
\end{proof}

\subsection{Why the functional-equation/conjugation pairing produces the block}

For completeness we isolate the algebra behind \eqref{ai:eq:Q}.  Weil's
Hermitian zero-side form pairs the evaluation at $\rho$ with the evaluation
at the horizontal partner under functional-equation plus conjugation symmetry
\[
 \rho^\sharp=1-\overline\rho.
\]
For an off-line pair of equal multiplicity $m_\rho$, its two contributions
combine as
\[
 m_\rho
 \left(
 \ell_\rho(c)\overline{\ell_{\rho^\sharp}(c)}
 +
 \ell_{\rho^\sharp}(c)\overline{\ell_\rho(c)}
 \right)
 =
 2m_\rho\operatorname{Re}
 \left(
 \ell_\rho(c)
 \overline{\ell_{\rho^\sharp}(c)}
 \right),
\]
which is exactly the hyperbolic block of Lemma~\ref{ai:lem:block}.  On the
critical line, $\rho=\rho^\sharp$, so the same expression collapses to a
positive square (up to the normalization convention already absorbed into
the zero-side form).

Thus the signature claim is a direct algebraic consequence of functional
equation symmetry once the zero-side Hermitian form has been written in
evaluation coordinates.

\subsection{Adversarial checks}

Several possible pathologies do not change the result.

\begin{itemize}
\item \textbf{Near-line pairs.}
As $\operatorname{Re}\rho\to1/2$, the two evaluation functionals may become
nearly dependent, but pull-back can only \emph{decrease} positive index.
No uniform linear independence is required.

\item \textbf{Deep off-line pairs.}
The eigenvalues $\pm m_\rho$ of the abstract pair block are independent of
the depth $|\operatorname{Re}\rho-1/2|$.  The evaluation map may distort or
annihilate directions, but cannot increase $n_+$.

\item \textbf{Coincident ordinates.}
Different zero points may have the same ordinate.  This creates dependence
among evaluation functionals, which again can only lower the pull-back
positive index.

\item \textbf{Multiplicity.}
A multiple critical-line point contributes one positive coordinate block
$m_\rho|z|^2$, hence one positive direction, regardless of $m_\rho$.
Multiplicity changes the weight but not its positive index.

\item \textbf{Noninjective compression.}
No injectivity of the full evaluation map is assumed; this is precisely why
Lemma~\ref{ai:lem:pullback} is formulated for arbitrary linear maps.
\end{itemize}



\section{Independent proof of Input II: exterior-zero tail}\label{app:inputII}
This appendix proves the trace-class tail estimate and its effect on the
rank--trace functional.  The exterior matrix is grouped into Hermitian
critical-line squares and off-line pair terms $uv^*+vu^*$, so that
$G=A+E$ holds entrywise under the same zero-side convention.
\subsection{Setup}

Let
\[
 \ell=\log\frac{T}{2\pi},\qquad L=\ell,\qquad
 X=e^L=\frac{T}{2\pi},\qquad D_0=T^{1/2},
\]
and let
\[
 I=(T,2T],\qquad I'=(T-D_0,\,2T+D_0].
\]
Let the critical sampling grid be
\[
 \tau_k=T+\frac{2\pi k}{L},\qquad 0\le k<d,
 \qquad d=\left\lfloor\frac{LT}{2\pi}\right\rfloor.
\]
The tapered Montgomery--Taylor window $\phi$ has support in
$[-L/2,L/2]$, satisfies $0\le\phi\le1$, and for fixed taper width has
\[
 C_1:=\|\phi''\|_1=O(1).
\]
Put
\[
 a=\frac1L\int_{\mathbb R}\phi(u)^2\,du.
\]
For this window, $a\to K(0)>0$; hence $a\asymp1$.

Use the globally fixed test family \eqref{eq:test-family}. Thus
\[
 G_{jk}=W(f_j,f_k)
\]
is exactly the common matrix \eqref{eq:G-common}.
Using the zero-side expansion of the normalized Weil formula, define
$A_{jk}$ by restricting that zero sum to zeros whose ordinary ordinates lie
in $I'$, grouping every off-line zero together with its horizontal partner
$\rho^\sharp=1-\overline\rho$.  Define $E_{jk}$ by the complementary
exterior zero sum.  Since the inside and outside sets form a disjoint
partition invariant under $\rho\mapsto\rho^\sharp$, one has entrywise
\[
 \boxed{G=A+E.}
\]
Equivalently,
\[
 \widehat A=\widehat G-\widehat E,
 \qquad
 \widehat G=\frac{G}{aL^2},\quad
 \widehat A=\frac{A}{aL^2},\quad
 \widehat E=\frac{E}{aL^2}.
\]

\subsection{Off-axis Fourier decay}

We use the standard Paley--Wiener integration-by-parts bound.

\begin{lemma}[Complex-shift decay]\label{aii:lem:complexdecay}
Let $\phi\in C_c^2([-L/2,L/2])$.  For real $r\ne0$ and
$|y|<1/2$,
\[
 |\widehat\phi(r-iy)|
 \le
 e^{L|y|/2}\,\|\phi''\|_1\,|r-iy|^{-2}
 \le
 X^{1/4}C_1\,|r|^{-2}.
\]
\end{lemma}

\begin{proof}
Twice integrating by parts in
\[
 \widehat\phi(r-iy)
 =\int_{-L/2}^{L/2}
 \phi(u)e^{-iru-yu}\,du
\]
and using the compactly supported $C^2$ taper gives
\[
 \widehat\phi(r-iy)
 =
 -(r-iy)^{-2}
 \int\phi''(u)e^{-i(r-iy)u}\,du.
\]
Since $|e^{-yu}|\le e^{L|y|/2}$ on the support,
\[
 |\widehat\phi(r-iy)|
 \le e^{L|y|/2}C_1|r-iy|^{-2}.
\]
Now $|r-iy|\ge|r|$ and, because $|y|<1/2$,
\[
 e^{L|y|/2}\le e^{L/4}=X^{1/4}.
\]
\end{proof}

\subsection{One external zero}

Let $\rho=\beta+i\gamma$ be a distinct zero with $\gamma\notin I'$ and put
\[
 y=\beta-\frac12,\qquad |y|<\frac12.
\]
Define its sampled vector
\[
 u_\rho=
 \bigl(\widehat\phi((\gamma-\tau_k)-iy)\bigr)_{0\le k<d}.
\]
Let
\[
 D=\operatorname{dist}(\gamma,I)\ge D_0.
\]

\begin{lemma}[Tail vector bound]\label{aii:lem:vector}
Uniformly for every such zero,
\[
 \|u_\rho\|_2^2
 \ll C_1^2X^{1/2}L D^{-3}.
\]
\end{lemma}

\begin{proof}
Lemma~\ref{aii:lem:complexdecay} gives
\[
 \|u_\rho\|_2^2
 \le
 C_1^2X^{1/2}
 \sum_{0\le k<d}|\gamma-\tau_k|^{-4}.
\]
The grid spacing is $h=2\pi/L$.  Comparing the monotone sum with an integral,
\[
 \sum_{0\le k<d}|\gamma-\tau_k|^{-4}
 \ll
 D^{-4}+\frac1h\int_D^\infty t^{-4}\,dt
 \ll
 D^{-4}+L D^{-3}
 \ll L D^{-3},
\]
because $D\ge1$ and $L\ge2$ for large $T$.
\end{proof}

\subsection{Summing over zeros outside the enlarged interval}

We use only the classical estimate
\begin{equation}\label{aii:eq:localzeros}
 N(t+1)-N(t)\le A_0\log(t+3)
\end{equation}
for some absolute $A_0$ and all $t\ge0$.

\begin{lemma}[Weighted exterior zero count]\label{aii:lem:zerotail}
Counting zeros with multiplicity,
\[
 \sum_{\gamma\notin I'}m_\rho\,
 \operatorname{dist}(\gamma,I)^{-3}
 \ll
 \frac{\log T}{D_0^2}.
\]
\end{lemma}

\begin{proof}
For the upper tail, group zeros into intervals
\[
 2T+D_0+j<\gamma\le2T+D_0+j+1,\qquad j\ge0.
\]
By \eqref{aii:eq:localzeros}, the contribution is
\[
 \ll
 \sum_{j\ge0}
 \frac{\log(2T+D_0+j+4)}{(D_0+j)^3}.
\]
Splitting at $j=T$ gives
\[
 \ll \frac{\log T}{D_0^2}.
\]
The lower positive tail $0<\gamma<T-D_0$ is identical.  Zeros with
$\gamma\le0$ are at distance $\gg T$ from $I$ and contribute
$O(T^{-2}\log T)$, which is smaller.  Adding the three ranges proves the
claim.
\end{proof}

\subsection{Trace-norm tail bound}

The exterior zero-side matrix must be grouped according to the
horizontal symmetry involution.  If $u_\rho$ is the sampled evaluation column,
then
\[
 E=
 \sum_{\substack{\rho\ {\rm exterior}\\ \rho=\rho^\sharp}}
 m_\rho u_\rho u_\rho^*
 +
 \sum_{\substack{\{\rho,\rho^\sharp\}\ {\rm exterior}\\
                  \rho\ne\rho^\sharp}}
 m_\rho\left(
 u_\rho u_{\rho^\sharp}^*
 +
 u_{\rho^\sharp}u_\rho^*
 \right).
\]
This is Hermitian term-by-term after pairing.  Moreover
\[
 \|uv^*+vu^*\|_1
 \le2\|u\|_2\|v\|_2
 \le\|u\|_2^2+\|v\|_2^2.
\]
Hence the scalar tail-vector estimate may be summed exactly as before.

\begin{proposition}[Independent tail estimate]\label{aii:prop:tail}
One has
\[
 \|E\|_1
 \ll
 C_1^2X^{1/2}L\frac{\log T}{D_0^2}.
\]
Consequently
\[
 \|\widehat E\|_1
 =
 \frac{\|E\|_1}{aL^2}
 \ll
 \frac{X^{1/2}\log T}{LD_0^2}.
\]
At $\lambda=1$, $L=\ell$, $X=T/(2\pi)$ and $D_0=T^{1/2}$, hence
\[
 \boxed{\|\widehat E\|_1\ll T^{-1/2}}.
\]
\end{proposition}

\begin{proof}
For critical-line exterior zeros use
$\|u_\rho u_\rho^*\|_1=\|u_\rho\|_2^2$.
For each off-line pair use the preceding inequality.  Both members have the
same ordinate and opposite horizontal displacement, so
Lemma~\ref{aii:lem:vector} applies to each.  Summing with multiplicity and using
Lemma~\ref{aii:lem:zerotail} gives the stated unnormalised bound.  Division by
$aL^2$, with $a\asymp1$, gives the normalized estimate.
\end{proof}

\subsection{The rank--trace functional under a trace-class perturbation}

Define
\[
 \mathcal R(H)=4\tr H-\|H\|_{\Frob}^2
\]
for a finite Hermitian matrix $H$.

\begin{lemma}[Perturbation inequality]\label{aii:lem:pert}
If $A=G-E$ are Hermitian matrices and $\varepsilon=\|E\|_1$, then
\begin{equation}\label{aii:eq:pert}
 \mathcal R(A)
 \ge
 \mathcal R(G)
 -
 4\varepsilon
 -
 2\|G\|_{\Frob}\varepsilon
 -
 \varepsilon^2.
\end{equation}
\end{lemma}

\begin{proof}
Since $|\tr E|\le\|E\|_1=\varepsilon$,
\[
 4\tr A\ge4\tr G-4\varepsilon.
\]
Also
\[
 \|A\|_{\Frob}
 \le
 \|G\|_{\Frob}+\|E\|_{\Frob}
 \le
 \|G\|_{\Frob}+\|E\|_1
 =
 \|G\|_{\Frob}+\varepsilon.
\]
Squaring and combining the two inequalities proves \eqref{aii:eq:pert}.
\end{proof}

\begin{theorem}[Independent Input II]\label{aii:thm:II}
For the optimized bandwidth $\lambda=1$, the prime-side moment estimate of Appendix~\ref{app:inputIII} gives $\|\widehat G\|_{\Frob}=O(N(T,2T)^{1/2})$.  Then
\[
 \boxed{
 4\tr\widehat A-\|\widehat A\|_{\Frob}^2
 \ge
 4\tr\widehat G-\|\widehat G\|_{\Frob}^2
 -o(N(T,2T)).
 }
\]
\end{theorem}

\begin{proof}
Apply Lemma~\ref{aii:lem:pert} to the normalized matrices with
\[
 \varepsilon=\|\widehat E\|_1\ll T^{-1/2}
\]
from Proposition~\ref{aii:prop:tail}.  Since
$N(T,2T)\asymp T\log T$,
\[
 \varepsilon=o(N),
 \qquad
 \varepsilon^2=o(N),
\]
and using the Appendix~\ref{app:inputIII} bound,
\[
 \|\widehat G\|_{\Frob}\varepsilon
 =
 O(N^{1/2}T^{-1/2})
 =
 o(N).
\]
Substitution in \eqref{aii:eq:pert} proves the result.
\end{proof}

\begin{remark}[Package closure]
Appendix~\ref{app:inputIII} establishes the prime-side second moment without using the present tail comparison.  Thus using $\|\widehat G\|_{\Frob}=O(N^{1/2})$ here is logically acyclic.
\end{remark}



\section{Independent proof of Input III: prime-side moments}\label{app:inputIII}
This appendix fixes the Fourier normalization of Weil's explicit formula~\cite{Weil52,Titchmarsh86},
derives the spectral density, controls integrated finite-grid end effects,
and evaluates the prime diagonal and both off-diagonal frequency types.
The weighted Montgomery--Vaughan inequality~\cite{MV74} is used only with an unspecified
absolute constant.
\subsection{Classical input and notation}

Put
\[
 \ell=\log\frac{T}{2\pi},\qquad L=\ell,\qquad
 X=e^L=\frac{T}{2\pi},\qquad I=[T,2T],
\]
and define
\begin{equation}\label{aiii:eq:ell1}
 \ell_1:=\ell+2\log2-1.
\end{equation}
Thus the Riemann--von Mangoldt formula reads
\[
 N(T,2T)=\frac{T\ell_1}{2\pi}+O(\ell).
\]
Let $\phi=\phi_L\in C_c^2(\R)$ be real, even, supported exactly on
$[-L/2,L/2]$, with a fixed-width endpoint taper and uniform derivative
bounds.  Define
\[
 \widehat\phi(r)=\int_\R \phi(u)e^{-iru}\,du,
 \qquad
 \Phi(r)=\int_\R \phi(u)^2e^{-iru}\,du,
\]
\[
 a=\frac1L\int\phi^2,\qquad
 b=\frac1L\int\phi^4,
 \qquad g=\phi^2*\phi^2.
\]
For the optimized family, $a$ and $b$ remain bounded above and below by
positive constants.

We use the classical Weil explicit formula in the following fixed
normalization. With
\[
 \widehat q(z)=\int_\R q(u)e^{izu}\,du,\qquad
 q(u)=\frac1{2\pi}\int_\R\widehat q(t)e^{-itu}\,dt,
\]
the external theorem used here is
\begin{align}
 \sum_\rho m_\rho\widehat q(z_\rho)
 &=
 \widehat q(i/2)+\widehat q(-i/2)
 -\sum_{n\ge1}\frac{\Lambda(n)}{\sqrt n}
 \bigl(q(\log n)+q(-\log n)\bigr)\nonumber\\
 &\quad+
 \frac1{2\pi}\int_\R\widehat q(\tau)
 \left[
 \Re\frac{\Gamma'}{\Gamma}
 \left(\frac14+\frac{i\tau}{2}\right)-\log\pi
 \right]d\tau .
 \label{aiii:eq:weil-classical}
\end{align}
Here $z_\rho=-i(\rho-\tfrac12)$ and zeros are counted with multiplicity.
This is the standard Guinand--Weil explicit formula in the present Fourier
normalization; see Weil~\cite{Weil52} and, for a modern formulation with
explicit strip hypotheses, Iwaniec--Kowalski~\cite[Theorem~5.12]{IK04}.
We apply it to
$q=f*\widetilde g$, $\widetilde g(u)=\overline{g(-u)}$.
Since $f,g\in C_c^2([-L/2,L/2])$, one has
$q\in C_c^2([-L,L])$.  Moreover
\[
 \widehat q(z)
 =
 \mathcal F_+f(z)\,
 \overline{\mathcal F_+g(\overline z)},
\]
and two integrations by parts in each factor give, uniformly on every fixed
strip $|\Im z|\le\tfrac12+\varepsilon$,
\[
 |\widehat q(z)|\ll_{\varepsilon,L,f,g}(1+|z|)^{-4}.
\]
Thus the standard strip-decay hypotheses of the explicit formula are
satisfied.  The support condition also implies that only
$n\le X=e^L$ occur in the prime sum. Fourier inversion then gives the
following proposition.

\begin{proposition}[Normalized Weil spectral density]\label{aiii:lem:spectral}
The compressed Weil form has density
\[
 \nu_X(\tau)=\mu(\tau)+\Pi_X(\tau)+P_X(\tau),
\]
where
\begin{align}
 \mu(\tau)
 &=\frac1{2\pi}\Re\frac{\Gamma'}{\Gamma}
 \left(\frac14+\frac{i\tau}{2}\right)
 -\frac{\log\pi}{2\pi},\\
 P_X(\tau)
 &=-\frac1\pi\sum_{n\le X}\frac{\Lambda(n)}{\sqrt n}
 \cos(\tau\log n),\\
 \Pi_X(\tau)
 &=\frac{1}{2\pi(1/4+\tau^2)}
 +\frac1\pi\Re\frac{X^{1/2+i\tau}-1}{1/2+i\tau}.
\end{align}
More precisely,
\[
 W(f,g)=\int_\R
 \mathcal F_+f(\tau)\overline{\mathcal F_+g(\tau)}
 \nu_X(\tau)\,d\tau.
\]
\end{proposition}

\begin{proof}
The gamma contribution in the classical explicit formula is the integral
against $\mu$.  For the prime term,
\[
 q(y)+q(-y)
 =\frac1\pi\int_\R
 \mathcal F_+f(\tau)\overline{\mathcal F_+g(\tau)}
 \cos(\tau y)\,d\tau,
\]
and $q(\pm\log n)=0$ for $n>X$.
For the pole terms,
\[
 H(i/2)+H(-i/2)
 =\int_{\R}q(u)(e^{-u/2}+e^{u/2})\,du.
\]
Since $\operatorname{supp}q\subset[-L,L]$, this equals
\[
 \int_{\R}q(u)e^{-|u|/2}\,du
 +
 \int_{-L}^{L}q(u)e^{|u|/2}\,du.
\]
With $s=1/2+i\tau$,
\[
 \int_\R e^{-|u|/2}e^{-i\tau u}du=(1/4+\tau^2)^{-1},
\]
while
\[
 \int_{-L}^Le^{|u|/2}e^{-i\tau u}du
 =
 2\Re\frac{X^{\bar s}-1}{\bar s}
 =
 2\Re\frac{X^s-1}{s}.
\]
Fourier inversion therefore gives exactly the density $\Pi_X$.  This proves the spectral identity with all signs and
$2\pi$ factors fixed.
\end{proof}

Stirling gives, uniformly for $\tau\in I$,
\[
 \mu(\tau)=\frac1{2\pi}\log\frac{\tau}{2\pi}+O(T^{-2}),
 \qquad \mu'(\tau)=O(T^{-1}).
\]
Also
\[
 |P_X(\tau)|\ll\sqrt X,\qquad
 |\Pi_X(\tau)|\ll\frac{\sqrt X}{1+|\tau|},
\]
and Riemann--von Mangoldt yields
\[
 \int_T^{2T}\mu(\tau)d\tau=N(T,2T)+O(\ell).
\]

Let
\[
 h=\frac{2\pi}{L},\qquad
 d=\left\lfloor\frac{T}{h}\right\rfloor,
 \qquad \tau_k=T+kh\quad(0\le k<d).
\]
By \eqref{eq:G-spectral-common}, the same common matrix $G$ satisfies
\[
 G_{jk}=\int_\R
 \widehat\phi(\tau-\tau_j)\widehat\phi(\tau-\tau_k)
 \nu_X(\tau)\,d\tau,
\]
and set
\[
 \widetilde G:=\frac{G}{L},\qquad
 \widehat G:=\frac{\widetilde G}{aL}=\frac{G}{aL^2}.
\]
This distinction between $G$, $\widetilde G$ and $\widehat G$ is maintained
throughout.

\subsection{Fourier and prime-power estimates}

The fixed-width $C^2$ taper gives
\begin{equation}\label{aiii:eq:psi}
 \max\{|\widehat\phi(r)|,|\Phi(r)|\}
 \le \psi(r):=
 \min\left\{L,\frac{C_1}{|r|},\frac{C_2}{r^2}\right\},
\end{equation}
with constants independent of $L$.  Consequently
\begin{equation}\label{aiii:eq:psimom}
 \Psi_0:=\int_0^\infty\psi(r)\,dr\ll\log L,
 \qquad
 \int_\R \psi(r)^2|r|\,dr\ll\log L,
 \qquad
 \int_\R\psi(r)^2\,dr\ll L.
\end{equation}
Plancherel and Fourier inversion give
\begin{equation}\label{aiii:eq:Phi}
 \int_\R\Phi(x)^2\,dx=2\pi bL,
 \qquad
 \int_\R\Phi(x)^2e^{ixy}\,dx=2\pi g(y).
\end{equation}

We also use the classical prime-power estimates
\begin{equation}\label{aiii:eq:primepower}
 \sum_{n\le x}\frac{\Lambda(n)}{\sqrt n}\ll\sqrt x,
 \qquad
 \sum_{n\le x}\Lambda(n)^2\ll x\log x,
\end{equation}
and
\begin{equation}\label{aiii:eq:primepower2}
 \sum_{n\le x}\frac{\Lambda(n)^2}{n}
 =\frac12\log^2x+O(\log x).
\end{equation}
The asymptotic \eqref{aiii:eq:primepower2} follows from the Prime Number
Theorem by standard partial summation; the upper bounds in
\eqref{aiii:eq:primepower} require only Chebyshev-strength estimates.
We use Titchmarsh~\cite{Titchmarsh86} as a standard reference.

\subsection{First trace moment}

Put
\[
 S_d(\tau):=\sum_{0\le k<d}|\widehat\phi(\tau-\tau_k)|^2,
 \qquad
 S_\infty(\tau):=\sum_{k\in\mathbb Z}
 |\widehat\phi(\tau-\tau_k)|^2.
\]
For the infinite translated critical grid, Poisson summation gives
\begin{equation}\label{aiii:eq:poisson-trace}
 S_\infty(\tau)=aL^2
 \qquad(\tau\in\mathbb R).
\end{equation}
We require two distinct end-effect estimates.  The first controls the lattice
points omitted from the finite grid while $\tau$ remains in $I$:
\[
 R_{\rm in}:=\int_I\bigl(S_\infty(\tau)-S_d(\tau)\bigr)
 |\nu_X(\tau)|\,d\tau.
\]
On $I$ one has
\[
 |\nu_X(\tau)|\ll B,\qquad B:=\ell+\sqrt X,
\]
by Stirling and the displayed bounds for $P_X,\Pi_X$.  Using
\eqref{aiii:eq:psi}, the left omitted tail is bounded by
\[
 \sum_{j\ge1}\int_{jh}^{\infty}\psi(r)^2\,dr,
\]
and the right omitted tail is identical up to $O(1)$ endpoint indices.
Hence, by Tonelli and \eqref{aiii:eq:psimom},
\begin{equation}\label{aiii:eq:Rin}
 R_{\rm in}\ll
 B\left(
 \frac1h\int_0^\infty r\psi(r)^2\,dr
 +\int_0^\infty\psi(r)^2\,dr
 \right)
 \ll BL\log L.
\end{equation}

The second estimate, absent from a naive truncation to $I$, controls the
actual finite-grid mass outside $I$:
\[
 R_{\rm out}:=\int_{\mathbb R\setminus I}
 S_d(\tau)|\nu_X(\tau)|\,d\tau.
\]
For the prime and pole pieces the global bounds
\[
 |P_X(\tau)|\ll\sqrt X,\qquad |\Pi_X(\tau)|\ll 1+\sqrt X
\]
give, by the same endpoint-tail summation,
\begin{equation}\label{aiii:eq:Rout-pp}
 \int_{\mathbb R\setminus I}S_d(\tau)
 (|P_X(\tau)|+|\Pi_X(\tau)|)\,d\tau
 \ll \sqrt X\,L\log L.
\end{equation}
For the archimedean term, Stirling gives
$|\mu(\tau)|\ll\ell$ on $|\tau|\le3T$.  Thus the part of
$\mathbb R\setminus I$ with $|\tau|\le3T$ is $O(\ell L\log L)$ by the same
argument.  If $|\tau|>3T$, then every $\tau_k\in[T,2T+O(h)]$ satisfies
$|\tau-\tau_k|\asymp|\tau|$, and \eqref{aiii:eq:psi} gives
\[
 S_d(\tau)\ll d\,|\tau|^{-4}\ll TL\,|\tau|^{-4},
 \qquad
 |\mu(\tau)|\ll\log(2+|\tau|).
\]
Consequently
\[
 \int_{|\tau|>3T}S_d(\tau)|\mu(\tau)|\,d\tau
 \ll TL\int_{3T}^\infty\frac{\log(2+t)}{t^4}\,dt
 \ll \frac{L\ell}{T^2}.
\]
Combining these estimates,
\begin{equation}\label{aiii:eq:Rout}
 R_{\rm out}\ll(\ell+\sqrt X)L\log L+\frac{L\ell}{T^2}
 =o(L^2N)
\end{equation}
for $L\asymp\ell$, $X\asymp T$, since $N\asymp T\ell$.

Now the spectral representation gives the exact decomposition
\[
 \tr G
 =
 \int_I S_d(\tau)\nu_X(\tau)\,d\tau
 +
 \int_{\mathbb R\setminus I}S_d(\tau)\nu_X(\tau)\,d\tau.
\]
Using \eqref{aiii:eq:poisson-trace}, \eqref{aiii:eq:Rin}, and
\eqref{aiii:eq:Rout},
\begin{equation}\label{aiii:eq:trace-reduction}
 \tr G
 =
 aL^2\int_I\nu_X(\tau)\,d\tau+o(L^2N).
\end{equation}
By Riemann--von Mangoldt and Stirling,
\[
 \int_I\mu(\tau)\,d\tau=N+O(\ell).
\]
Termwise integration and \eqref{aiii:eq:primepower} give
\[
 \left|\int_I P_X(\tau)\,d\tau\right|
 \ll\sum_{n\le X}\frac{\Lambda(n)}{\sqrt n\log n}
 \ll\sqrt X=o(N),
\]
while the explicit formula for $\Pi_X$ gives
$\int_I\Pi_X(\tau)\,d\tau=o(N)$.  Therefore
\begin{equation}\label{aiii:eq:first}
 \tr\widetilde G=aLN(1+o(1)),
 \qquad
 \boxed{\tr\widehat G=N(1+o(1)).}
\end{equation}

\subsection{Integrated finite-grid end effects for the second moment}

Define
\[
 K(\tau,\tau')=
 \sum_{0\le k<d}\widehat\phi(\tau-\tau_k)
 \widehat\phi(\tau'-\tau_k),
\]
\[
 K_\infty(\tau,\tau')=
 \sum_{k\in\mathbb Z}\widehat\phi(\tau-\tau_k)
 \widehat\phi(\tau'-\tau_k)=L\Phi(\tau-\tau'),
\]
and $K_{\rm out}=K_\infty-K$.  Since
$\sum_{k\in\mathbb Z}|\widehat\phi(\tau-\tau_k)|^2=aL^2$,
Cauchy--Schwarz gives
\begin{equation}\label{aiii:eq:Kbound}
 |K|,|K_\infty|\le aL^2,
 \qquad
 |K_{\rm out}(\tau,\tau')|
 \le\sum_{k\notin[0,d)}\psi(\tau-\tau_k)\psi(\tau'-\tau_k).
\end{equation}
Put
\[
 B=C(\ell+\sqrt X),
\]
so $|\nu_X(\tau)|\le B$ for $|\tau|\le4T$ and
$B^2\ll \ell^2+X$.

\begin{lemma}[Global density bound used in the tails]\label{aiii:lem:nuglobal}
For $|\tau|\le4T$,
\[
 |\nu_X(\tau)|\ll \ell+\sqrt X=:B.
\]
If $\tau_k\in[T,2T]$ and $r=|\tau-\tau_k|>2T$, then
\[
 |\nu_X(\tau)|\ll B+\log(r/T).
\]
\end{lemma}

\begin{proof}
The prime term satisfies $|P_X|\ll\sqrt X$ by
$\sum_{n\le X}\Lambda(n)/\sqrt n\ll\sqrt X$.
The pole density satisfies
$|\Pi_X(\tau)|\ll \sqrt X/(1+|\tau|)+1/(1+\tau^2)$.
Stirling gives $\mu(\tau)\ll1+\log(2+|\tau|)$.
This proves the first estimate on $|\tau|\le4T$.
For $r>2T$ and $\tau_k\in[T,2T]$, one has
$|\tau|\le r+2T<2r$ and hence
$\log(2+|\tau|)\ll\ell+\log(r/T)$; the prime term remains bounded by
$O(\sqrt X)$ and the pole term is smaller.
\end{proof}

\begin{proposition}[Integrated end effects]\label{aiii:prop:end}
For $I=(T,2T]$ and
\[
 M=\iint_{I\times I}
 \Phi(\tau-\tau')^2\nu_X(\tau)\nu_X(\tau')\,d\tau d\tau',
\]
one has
\begin{equation}\label{aiii:eq:end-final}
 \tr\widetilde G^2
 =
 M+
 O\!\left(
 L\ell\log L\,(\ell^2+X)
 \right),
 \qquad \widetilde G=G/L.
\end{equation}
\end{proposition}

\begin{proof}
Expanding $\tr G^2$ gives
\[
 L^2\tr\widetilde G^2
 =
 \iint_{\R^2}|K(\tau,\tau')|^2
 \nu_X(\tau)\nu_X(\tau')\,d\tau d\tau'.
\]
On $I\times I$, $|K_\infty|^2=L^2\Phi^2$.  Write the difference from
$L^2M$ as $E_1+E_2$, where
\[
 E_1=
 \iint_{I\times I}(|K|^2-|K_\infty|^2)\nu\nu',
\qquad
 E_2=
 \iint_{(I\times I)^c}|K|^2\nu\nu'.
\]

For $E_1$,
\[
 \bigl||K|^2-|K_\infty|^2\bigr|
 \le |K-K_\infty|(|K|+|K_\infty|)
 =|K_{\rm out}|(|K|+|K_\infty|)
 \ll L^2|K_{\rm out}|.
\]
Hence
\[
 |E_1|
 \ll
 L^2B^2
 \sum_{k\notin[0,d)}
 \left(\int_I\psi(\tau-\tau_k)\,d\tau\right)^2.
\]
For omitted grid points at distance $\Delta>0$ from $I$,
\[
 \int_I\psi(\tau-\tau_k)d\tau
 \le\int_\Delta^\infty\psi(r)dr
 \le\min\left\{\Psi_0,\frac{C}{\Delta}\right\}.
\]
At the left edge, omitted points are $k=-j$, $j\ge1$, and have
distance $jh$ from $I$.  At the right edge the two exceptional omitted
indices $k=d,d+1$ may lie in, or within one mesh step of, the endpoint
$2T$ because $d=\lfloor T/h\rfloor$; for them we use the crude bound
$2\Psi_0$.  Every further right-hand omitted index has distance at least
$jh$, $j\ge1$.  Therefore
\begin{align*}
 S&:=
 \sum_{k\notin[0,d)}
 \left(\int_I\psi(\tau-\tau_k)d\tau\right)^2\\
 &\ll
 \Psi_0^2+
 \sum_{j\ge1}
 \min\left\{\Psi_0,\frac{C}{jh}\right\}^2.
\end{align*}
Let $J=\lceil C/(h\Psi_0)\rceil$.  Then
\[
 \sum_{j\le J}\Psi_0^2
 \ll\frac{\Psi_0}{h}+\Psi_0^2,
\]
whereas
\[
 \sum_{j>J}\frac{C^2}{j^2h^2}
 \ll\frac{C^2}{h^2J}
 \ll\frac{\Psi_0}{h}.
\]
Since $h=2\pi/L$ and $\Psi_0=O(\log L)$,
\[
 S=O(L\log L).
\]
Thus
\begin{equation}\label{aiii:eq:E1final}
 |E_1|\ll L^3B^2\log L.
\end{equation}

For $E_2$, symmetry lets us place the first variable outside $I$ and multiply
by two.  Since $|K|^2\ll L^2|K|$ and
$|K|\le\sum_{0\le k<d}\psi(\tau-\tau_k)\psi(\tau'-\tau_k)$,
\[
 |E_2|
 \ll
 L^2\sum_{0\le k<d}A_kC_k,
\]
where
\[
 A_k=\int_{\tau\notin I}\psi(\tau-\tau_k)|\nu_X(\tau)|\,d\tau,
\quad
 C_k=\int_{\R}\psi(\tau'-\tau_k)|\nu_X(\tau')|\,d\tau'.
\]

We first bound $C_k$ uniformly.  On
$|\tau'-\tau_k|\le2T$ one has $|\tau'|\le4T$, hence
\[
 \int_{|\tau'-\tau_k|\le2T}
 \psi(\tau'-\tau_k)|\nu_X(\tau')|d\tau'
 \le2B\Psi_0.
\]
For $r=|\tau'-\tau_k|>2T$, one has
$|\tau'|\le2r$ and
\[
 |\nu_X(\tau')|\ll B+\log(r/T),\qquad
 \psi(r)\ll r^{-2}.
\]
Therefore the far part is
\[
 \ll
 \int_{2T}^{\infty}
 \frac{B+\log(r/T)}{r^2}\,dr
 \ll \frac{B+1}{T}
 \ll B,
\]
and hence
\begin{equation}\label{aiii:eq:Ck}
 C_k\ll B\Psi_0.
\end{equation}

Now sum $A_k$ before integrating.  Put
\[
 \sigma(\tau)=\sum_{0\le k<d}\psi(\tau-\tau_k),
 \qquad
 \Delta=\operatorname{dist}(\tau,I).
\]
The distances to grid points dominate
$\Delta,\Delta+h,\Delta+2h,\ldots$, so monotonicity yields
\begin{equation}\label{aiii:eq:sigma1}
 \sigma(\tau)
 \le
 \psi(\Delta)+\frac1h\int_\Delta^\infty\psi(r)dr
 \ll
 \psi(\Delta)+L\min\left\{\Psi_0,\frac{C}{\Delta}\right\}.
\end{equation}
Also, since $d\ll LT$ and $\psi(r)\ll r^{-2}$,
\begin{equation}\label{aiii:eq:sigma2}
 \sigma(\tau)\ll LT\,\Delta^{-2}.
\end{equation}

For $0<\Delta\le2T$, $|\tau|\le4T$, so $|\nu_X|\le B$.  Integrating
\eqref{aiii:eq:sigma1} on the two sides of $I$ gives
\begin{align*}
 \int_{0<\Delta\le2T}|\nu_X(\tau)|\sigma(\tau)d\tau
 &\ll
 B\left[
 \Psi_0+
 L\int_0^{2T}
 \min\left\{\Psi_0,\frac{C}{\Delta}\right\}d\Delta
 \right]\\
 &\ll
 B\left[\Psi_0+L+L\log(2T\Psi_0/C)\right]\\
 &\ll BL\ell.
\end{align*}
For $\Delta>2T$, use \eqref{aiii:eq:sigma2} and
$|\nu_X(\tau)|\ll B+\log(\Delta/T)$:
\begin{align*}
 \int_{\Delta>2T}|\nu_X(\tau)|\sigma(\tau)d\tau
 &\ll
 LT\int_{2T}^{\infty}
 \frac{B+\log(\Delta/T)}{\Delta^2}d\Delta\\
 &\ll LB.
\end{align*}
Consequently
\begin{equation}\label{aiii:eq:Asum}
 \sum_kA_k
 =
 \int_{\tau\notin I}|\nu_X(\tau)|\sigma(\tau)d\tau
 \ll BL\ell.
\end{equation}
Combining \eqref{aiii:eq:Ck} and \eqref{aiii:eq:Asum},
\begin{equation}\label{aiii:eq:E2final}
 |E_2|
 \ll
 L^2(B\Psi_0)(BL\ell)
 \ll
 L^3B^2\ell\log L.
\end{equation}
Equations \eqref{aiii:eq:E1final} and \eqref{aiii:eq:E2final} give
\[
 L^2\tr\widetilde G^2
 =
 L^2M+
 O(L^3B^2\ell\log L).
\]
Divide by $L^2$ and use
$B^2\ll\ell^2+X$ to obtain \eqref{aiii:eq:end-final}.
\end{proof}

At $L=\ell$ and $X\asymp T$, the error in \eqref{aiii:eq:end-final} is
$o(TL\ell^2)$.


\subsection{Evaluation of the continuous second moment}

For functions $u_1,u_2$ on $I$, write
\[
 \mathcal M[u_1,u_2]=
 \iint_{I\times I}\Phi(\tau-\tau')^2
 u_1(\tau)u_2(\tau')\,d\tau\,d\tau'.
\]
Then
\[
 M=\mathcal M[\mu,\mu]+\mathcal M[P_X,P_X]
 +2\mathcal M[\mu,P_X]+2\mathcal M[\mu,\Pi_X]
 +2\mathcal M[P_X,\Pi_X]+\mathcal M[\Pi_X,\Pi_X].
\]

\subsubsection{Archimedean term}

Using $\mu(\tau)=\mu(\tau')+O(|\tau-\tau'|/T)$,
\eqref{aiii:eq:psimom}, and \eqref{aiii:eq:Phi}, one obtains
\begin{equation}\label{aiii:eq:arch}
 \mathcal M[\mu,\mu]
 =2\pi bL\int_I\mu(\tau)^2\,d\tau
 +O(\ell^2\log L)
 =\frac{bLT\ell_1^2}{2\pi}+o(TL\ell^2).
\end{equation}
Indeed, the local replacement error is bounded by
$O(\ell\int\Phi(x)^2|x|dx)$ after integrating over $I$, and the truncation
of the full $x$-integral at the two interval ends costs
$O(\ell^2\int\Phi(x)^2\min(|x|,T)dx)$.

\subsubsection{Prime term: diagonal and two off-diagonal families}

Write
\[
 a_n=\frac{\Lambda(n)}{\sqrt n},\qquad y_n=\log n,
\]
so
\[
 P_X(\tau)=-\frac1{2\pi}\sum_{n\le X}a_n
 (e^{i\tau y_n}+e^{-i\tau y_n}).
\]
Substitute $x=\tau-\tau'$.  For $|x|<T$, the $\tau'$-interval is
$[T+x_-,2T-x_+]$, $x_\pm=\max(\pm x,0)$.  Expanding the four exponential
products splits $\mathcal M[P_X,P_X]$ into a quotient-frequency diagonal
$D$, a quotient-frequency off-diagonal $O_1$, and a product-frequency term
$O_2$.

For $n=m$, the inner interval has length $T-|x|$, hence Fourier inversion
and \eqref{aiii:eq:psimom} give
\begin{equation}\label{aiii:eq:D}
 D=\frac{T}{\pi}\sum_{n\le X}a_n^2g(y_n)
 +O(L^2\log L).
\end{equation}
For product frequencies,
\[
 \left|\int_{T+x_-}^{2T-x_+}
 e^{i\tau'\log(nm)}\,d\tau'\right|
 \le\frac2{\log(nm)}
 \le\frac2{\log4}.
\]
Returning to the expression before the $x$-integration and taking absolute
values therefore gives
\begin{align*}
 |O_2|
 &\le
 \frac{1}{2\pi^2}
 \sum_{n,m\le X}a_na_m
 \int_{-T}^{T}\Phi(x)^2
 \frac{2}{\log(nm)}\,dx\\
 &\le
 \frac{1}{\pi^2\log4}
 \left(\sum_{n\le X}a_n\right)^2
 \int_{\R}\Phi(x)^2\,dx.
\end{align*}
By the prime-power estimate
\[
 \sum_{n\le X}a_n
 =
 \sum_{n\le X}\frac{\Lambda(n)}{\sqrt n}
 \ll \sqrt X
\]
and Plancherel
\[
 \int_{\R}\Phi(x)^2\,dx=2\pi bL\ll L.
\]
Consequently
\begin{equation}\label{aiii:eq:O2}
 \boxed{O_2\ll XL.}
\end{equation}

For $n\ne m$, put $\vartheta=y_n-y_m$ and
\[
 \alpha_n^+=\int_0^T\Phi(x)^2e^{ixy_n}\,dx,\qquad
 \alpha_n^-=\int_{-T}^{0}\Phi(x)^2e^{ixy_n}\,dx.
\]
Then $|\alpha_n^\pm|\le\pi bL\ll L$.  Direct evaluation of the inner
$\tau'$-integral gives
\begin{align*}
 O_1
 &=
 \frac1{2\pi^2}\Re\sum_{n\ne m}
 \frac{a_na_m}{i(y_n-y_m)}
 \Bigg[
 \left(\frac nm\right)^{2iT}
 (\alpha_m^++\alpha_n^-)\\
 &\hspace{40mm}-
 \left(\frac nm\right)^{iT}
 (\alpha_n^++\alpha_m^-)
 \Bigg].
\end{align*}
Thus $O_1$ is the sum of four Hilbert forms with coefficient pairs
\[
(a_nn^{2iT},a_mm^{-2iT}\alpha_m^+),\quad
(a_nn^{2iT}\alpha_n^-,a_mm^{-2iT}),
\]
\[
(a_nn^{iT}\alpha_n^+,a_mm^{-iT}),\quad
(a_nn^{iT},a_mm^{-iT}\alpha_m^-).
\]
We use the weighted Montgomery--Vaughan Hilbert inequality
\[
 \left|\sum_{r\ne s}\frac{x_rz_s}{\lambda_r-\lambda_s}\right|
 \ll
 \left(\sum_r\frac{|x_r|^2}{\delta_r}\right)^{1/2}
 \left(\sum_r\frac{|z_r|^2}{\delta_r}\right)^{1/2},
\quad
\delta_r=\min_{s\ne r}|\lambda_r-\lambda_s|,
\]
with an absolute implied constant; no sharp constant is needed.
For $\lambda_n=\log n$,
$\delta_n^{-1}\ll n$, since adjacent integers already have logarithmic
spacing $\gg1/n$.  Hence
\[
 \sum_{n\le X}\frac{a_n^2}{\delta_n}
 \ll\sum_{n\le X}\Lambda(n)^2\ll XL,
\]
and a sequence carrying $\alpha_n^\pm$ has the corresponding weighted
square sum $\ll L^3X$.  Each of the four Hilbert forms is therefore
$O(L^2X)$, so
\begin{equation}\label{aiii:eq:O1}
 O_1\ll L^2X.
\end{equation}
Combining \eqref{aiii:eq:D}--\eqref{aiii:eq:O1},
\begin{equation}\label{aiii:eq:prime}
 \mathcal M[P_X,P_X]
 =\frac{T}{\pi}\sum_{n\le X}\frac{\Lambda(n)^2}{n}g(\log n)
 +O(L^2X).
\end{equation}

By Stieltjes partial summation from \eqref{aiii:eq:primepower2}, and using
$\|g'\|_\infty\le\|(\phi^2)'\|_1=O(1)$,
\begin{equation}\label{aiii:eq:partial}
 \sum_{n\le X}\frac{\Lambda(n)^2}{n}g(\log n)
 =\int_0^L y g(y)\,dy+O(L^2).
\end{equation}

\subsubsection{Mixed and pole terms}

\begin{proposition}[Secondary terms]\label{aiii:prop:secondary}
At $L=\ell$ and $X\asymp T$,
\[
 \mathcal M[\mu,P_X]\ll \ell\sqrt X,
 \qquad
 \mathcal M[\mu,\Pi_X]\ll \ell L\sqrt X,
\]
\[
 \mathcal M[P_X,\Pi_X]\ll LX,
 \qquad
 \mathcal M[\Pi_X,\Pi_X]\ll \frac{LX}{T}.
\]
All four are $o(TL\ell^2)$.
\end{proposition}

\begin{proof}
Define
\[
 m(\tau')=\int_I\Phi(\tau-\tau')^2\mu(\tau)\,d\tau.
\]
Then $\|m\|_\infty\ll\ell L$.  Differentiating under the integral and
integrating by parts in $\tau$ gives
\[
 \int_I|m'(\tau')|\,d\tau'\ll\ell L,
\]
using $\mu'=O(T^{-1})$ and \eqref{aiii:eq:Phi}.  Therefore, for
$y\ge\log2$,
\[
 \left|\int_I m(\tau')e^{iy\tau'}d\tau'\right|
 \ll\frac{\ell L}{y}.
\]
Substitution of the defining formula for $P_X$ and
$\sum_{n\le X}\Lambda(n)/(\sqrt n\log n)\ll\sqrt X/L$ yields
$\mathcal M[\mu,P_X]\ll\ell\sqrt X$.

For the remaining terms use on $I$
\[
 |\Pi_X|\ll\frac{\sqrt X}{T},\qquad
 0<\mu\ll\ell,\qquad |P_X|\ll\sqrt X,
\]
and
\[
 \sup_\tau\int_I\Phi(\tau-\tau')^2d\tau'\le2\pi bL\ll L.
\]
Thus
\[
 |\mathcal M[\mu,\Pi_X]|
 \ll T\cdot\ell\cdot\frac{\sqrt X}{T}\cdot L
 =\ell L\sqrt X,
\]
\[
 |\mathcal M[P_X,\Pi_X]|
 \ll T\cdot\sqrt X\cdot\frac{\sqrt X}{T}\cdot L=LX,
\]
and
\[
 |\mathcal M[\Pi_X,\Pi_X]|
 \ll T\cdot\frac{X}{T^2}\cdot L=\frac{LX}{T}.
\]
\end{proof}

\subsection{Scale-free limit}

Assume the fixed-width tapered family satisfies
\[
 \phi_L(Ls)^2=v(s)+o(1)
\]
in $L^1([-1/2,1/2])$, with $v$ bounded and supported on
$[-1/2,1/2]$.  Then
\[
 a\to A:=\int_{-1/2}^{1/2}v(s)\,ds,
 \qquad
 b\to B_v:=\int_{-1/2}^{1/2}v(s)^2\,ds.
\]
Writing
\[
 G_v(r)=\int v(s)v(s-r)\,ds,
\]
one has $g(Lr)=LG_v(r)+o(L)$ in $L^1$, and hence
\begin{equation}\label{aiii:eq:J}
 \frac{2}{L^3}\int_0^L yg(y)\,dy
 \to J_v:=\iint_{[-1/2,1/2]^2}|s-t|v(s)v(t)\,ds\,dt.
\end{equation}

Combining Proposition~\ref{aiii:prop:end}, \eqref{aiii:eq:arch}, \eqref{aiii:eq:prime},
\eqref{aiii:eq:partial}, Proposition~\ref{aiii:prop:secondary}, and
\eqref{aiii:eq:J},
\[
 \tr\widetilde G^2
 =\frac{TL}{2\pi}\left(B_v\ell_1^2+L^2J_v\right)
 +o(TL\ell^2).
\]
Since, by \eqref{aiii:eq:ell1},
\[
 N(T,2T)=\frac{T\ell_1}{2\pi}+O(\ell),
 \qquad \frac{L}{\ell_1}\to1,
\]
and $\widehat G=\widetilde G/(aL)$,
\begin{equation}\label{aiii:eq:momentgeneral}
 \boxed{
 \|\widehat G\|_{\Frob}^2
 =\left(\frac{B_v+J_v}{A^2}+o(1)\right)N(T,2T).
 }
\end{equation}
Thus
\[
 c_1(v)=\frac{A^2}{B_v+J_v}.
\]
Together with \eqref{aiii:eq:first}, this proves the two required prime-side
moments for every admissible profile.

\subsection{Closing the Montgomery--Taylor optimization}

On $H=L^2([-1/2,1/2])$ define the self-adjoint integral operator
\[
 (\mathcal Tv)(s)=\int_{-1/2}^{1/2}|s-t|v(t)\,dt.
\]
Then
\[
 c_1(v)=
 \frac{\langle1,v\rangle^2}
 {\langle v,(I+\mathcal T)v\rangle}.
\]

\begin{lemma}[Positivity and invertibility]\label{aiii:lem:positive}
$I+\mathcal T$ is strictly positive and invertible on $H$.
\end{lemma}

\begin{proof}
Schur's test gives
\[
 \|\mathcal T\|_{2\to2}
 \le\sup_{|s|\le1/2}\int_{-1/2}^{1/2}|s-t|\,dt
 =\sup_{|s|\le1/2}\left(s^2+\frac14\right)=\frac12.
\]
Hence
\[
 \langle v,(I+\mathcal T)v\rangle
 \ge(1-\|\mathcal T\|)\|v\|_2^2
 \ge\frac12\|v\|_2^2.
\]
\end{proof}

Cauchy--Schwarz in the positive inner product
$\langle f,g\rangle_A=\langle f,(I+\mathcal T)g\rangle$ now gives
\[
 c_1(v)\le\langle1,(I+\mathcal T)^{-1}1\rangle,
\]
with equality precisely when
\[
 v\propto(I+\mathcal T)^{-1}1.
\]
Thus a maximizer satisfies
\begin{equation}\label{aiii:eq:inteq}
 (I+\mathcal T)v=C
\end{equation}
for some constant $C$.
In the interior,
\[
 (\mathcal Tv)''=2v,
\]
so \eqref{aiii:eq:inteq} implies
\[
 v''+2v=0.
\]
By symmetry the maximizing solution is even, hence
\[
 v(s)=A_0\cos(\sqrt2s).
\]
Conversely, for $v_*(s)=\cos(\sqrt2s)$ the function
$F=v_*+\mathcal Tv_*$ satisfies $F''=0$ and is even, hence is constant;
therefore $v_*$ really satisfies \eqref{aiii:eq:inteq}, not merely its
differentiated equation.  Since $1/\sqrt2<\pi/2$, $v_*>0$ on the whole
interval, so the constraint $v\ge0$ is inactive.

Let $\theta=1/\sqrt2$.  Then
\[
 A=\int_{-1/2}^{1/2}v_*(s)\,ds=\sqrt2\sin\theta.
\]
Evaluating $C=(I+\mathcal T)v_*$ at $s=0$ gives
\[
 C=1+2\int_0^{1/2}t\cos(\sqrt2t)\,dt
 =\cos\theta+\frac{\sin\theta}{\sqrt2}.
\]
Since $\langle v_*,(I+\mathcal T)v_*\rangle=CA$,
\[
 c_1^*=\frac{A}{C}
 =\frac{2\tan(1/\sqrt2)}{\sqrt2+\tan(1/\sqrt2)},
\]
and therefore
\begin{equation}\label{aiii:eq:cstar}
 \boxed{
 \frac1{c_1^*}
 =\frac12+\frac1{\sqrt2}\cot\frac1{\sqrt2}.
 }
\end{equation}



\section{Independent proof of Input IV: Poisson--Gabor sampling}\label{app:inputIV}
This appendix proves the critical-lattice Poisson identity, the tapered
Fourier decay, finite-grid truncation, and the Montgomery--Taylor overlap
limit.  The window is real and even, so the real-axis transforms used below
are real and even; modulus squares may be retained throughout.
\subsection{Statement to be verified}

Let $L\to\infty$, put
\[
 h=\frac{2\pi}{L},\qquad \tau_k=kh,
\]
and let $\phi_L\in C_c^2(\mathbb R)$ be a real tapered window with exact support
$[-L/2,L/2]$.  Write
\[
 \widehat\phi_L(\xi)=\int_{\R}\phi_L(u)e^{-i\xi u}\,du
\]
and
\[
 a_L=\frac1L\int_{\R}\phi_L(u)^2\,du.
\]
For the Montgomery--Taylor choice, away from an $O(1)$ physical taper layer,
\[
 \phi_L(u)^2=\cos\!\left(\frac{\sqrt2\,u}{L}\right),
 \qquad |u|\le \frac L2-O(1).
\]
Define
\[
 K(x)=\int_{-1/2}^{1/2}\cos(\sqrt2\,t)\cos(2\pi xt)\,dt,
 \qquad k(x)=\frac{K(x)}{K(0)}.
\]

The required form of Input IV is the following.

\begin{theorem}[Independent Input IV]\label{aiv:thm:IV}
Fix $R>0$.  Let $\gamma,\gamma'$ be points whose normalized coordinates
$x=L\gamma/(2\pi)$ and $x'=L\gamma'/(2\pi)$ are at least $L^2$ grid cells
from both ends of the finite sampling interval, and suppose
$|x-x'|\le R$.  Then the normalized finite-grid Gram entry satisfies
\[
 \frac{1}{a_LL^2}
 \sum_{\tau_j\ {\rm in\ the\ finite\ grid}}
 \widehat\phi_L(\gamma-\tau_j)
 \overline{\widehat\phi_L(\gamma'-\tau_j)}
 =
 k(x-x')+o(1),
\]
uniformly for $|x-x'|\le R$.
\end{theorem}

\subsection{The infinite-grid identity}

\begin{lemma}[Critical-grid Poisson identity]\label{aiv:lem:poisson}
Assume $\phi\in L^2(\R)$ is compactly supported in an interval of length at
most $L$ and is sufficiently regular for Poisson summation.  Then
\[
 \sum_{k\in\mathbb Z}
 \widehat\phi(\tau-\tau_k)
 \overline{\widehat\phi(\tau'-\tau_k)}
 =
 L\int_{\R}|\phi(u)|^2e^{-i(\tau-\tau')u}\,du .
\]
\end{lemma}

\begin{proof}
Set
\[
 F_{\tau,\tau'}(\xi)
 =
 \widehat\phi(\tau-\xi)\,
 \overline{\widehat\phi(\tau'-\xi)}.
\]
With the Fourier convention above, a direct Fourier inversion calculation
gives
\[
 \widehat F_{\tau,\tau'}(s)
 =
 2\pi e^{-i\tau' s}
 \int_{\R}\phi(u)\overline{\phi(u-s)}
 e^{-i(\tau-\tau')u}\,du.
\]
Poisson summation on the lattice $h\mathbb Z$, $h=2\pi/L$, gives
\[
 \sum_{k\in\mathbb Z}F_{\tau,\tau'}(kh)
 =
 \frac1h\sum_{m\in\mathbb Z}
 \widehat F_{\tau,\tau'}(mL).
\]
Because $\phi$ is supported in an interval of length at most $L$, the
translated supports of $\phi(u)$ and $\phi(u-mL)$ are disjoint for
$m\ne0$, apart from endpoints of measure zero.  Thus only $m=0$ survives.
Since $1/h=L/(2\pi)$,
\[
 \sum_{k\in\mathbb Z}F_{\tau,\tau'}(kh)
 =
 L\int_{\R}|\phi(u)|^2e^{-i(\tau-\tau')u}\,du.
\]
\end{proof}

\begin{remark}
The exact support condition $\operatorname{supp}\phi_L=[-L/2,L/2]$ is the one used by the optimized fixed-width tapered family.  Hence the critical-density Poisson identity has no aliasing term: shifts by $mL$, $m\ne0$, overlap only on sets of measure zero.
\end{remark}

\subsection{Fourier decay}

\begin{lemma}[Uniform decay]\label{aiv:lem:decay}
Suppose $\phi_L\in C_c^2([-L/2,L/2])$ is uniformly bounded and its taper has fixed physical width with uniformly bounded first and second derivatives.  Then
\[
 |\widehat\phi_L(r)|
 \ll
 \min\left\{L,\frac1{|r|},\frac1{r^2}\right\},
\]
with an implied constant independent of $L$.
\end{lemma}

\begin{proof}
The $L^1$ estimate gives $|\widehat\phi_L(r)|\ll L$.
One integration by parts, using compact support and the bounded total
variation of $\phi_L$, gives $|\widehat\phi_L(r)|\ll |r|^{-1}$.
A second integration by parts on each smooth piece gives
$|\widehat\phi_L(r)|\ll |r|^{-2}$: the bulk derivatives of the
Montgomery--Taylor profile contribute $O(L^{-1})$ and $O(L^{-2})$ per unit
length, while the fixed-width taper contributes $O(1)$ total variation at
each derivative level.  All boundary terms vanish because the taper is compactly supported and $C^2$.
\end{proof}

\subsection{Finite-grid truncation}

\begin{lemma}[Pointwise tail]\label{aiv:lem:tail}
Fix $R>0$.  Suppose $\gamma,\gamma'$ are at least $L^2$ grid cells from a
finite-grid endpoint and $|L(\gamma-\gamma')/(2\pi)|\le R$.  Then the
contribution of the omitted half-grid beyond that endpoint is $O(L^{-2})$
before Gram normalization.
\end{lemma}

\begin{proof}
Let the omitted index be $n\ge L^2$ cells beyond the nearer endpoint.  Since
$h=2\pi/L$ and the normalized separation of $\gamma,\gamma'$ is bounded,
\[
 |\gamma-\tau_j|\asymp \frac nL,\qquad
 |\gamma'-\tau_j|\asymp \frac nL
\]
uniformly in the stated region.  Lemma~\ref{aiv:lem:decay} therefore gives
\[
 |\widehat\phi_L(\gamma-\tau_j)
   \widehat\phi_L(\gamma'-\tau_j)|
 \ll \frac{L^4}{n^4}.
\]
Consequently
\[
 \sum_{n\ge L^2}\frac{L^4}{n^4}
 \ll L^4\int_{L^2-1}^{\infty}t^{-4}\,dt
 =O(L^{-2}).
\]
The same estimate applies at the other endpoint.
\end{proof}

\subsection{Montgomery--Taylor overlap limit}

\begin{lemma}[Normalization and overlap]\label{aiv:lem:overlap}
For the fixed-width tapered Montgomery--Taylor family,
\[
 a_L=K(0)+o(1),
 \qquad
 K(0)=\int_{-1/2}^{1/2}\cos(\sqrt2 t)\,dt
 =\sqrt2\sin(1/\sqrt2)>0.
\]
Moreover, uniformly for $|x|\le R$,
\[
 \frac1{a_LL}
 \int_{\R}\phi_L(u)^2e^{-i(2\pi x/L)u}\,du
 =
 k(x)+o(1).
\]
\end{lemma}

\begin{proof}
Put $u=Lt$.  Outside a set of $t$-measure $O(L^{-1})$ arising from the
fixed-width taper,
\[
 \phi_L(Lt)^2=\cos(\sqrt2 t)
\]
on $[-1/2,1/2]$.  Uniform boundedness of the window makes the taper
contribution $O(L^{-1})$ after normalization.  Dominated convergence,
uniformly for $|x|\le R$, gives
\[
 a_L
 =
 \int_{-1/2}^{1/2}\cos(\sqrt2 t)\,dt+o(1)
 =
 K(0)+o(1)
\]
and
\[
 \frac1L\int\phi_L(u)^2e^{-i(2\pi x/L)u}\,du
 =
 \int_{-1/2}^{1/2}
 \cos(\sqrt2 t)e^{-i2\pi xt}\,dt+o(1).
\]
The odd imaginary part vanishes.  Dividing by $a_L=K(0)+o(1)$ proves the
claim.
\end{proof}

The overlap has the explicit form
\[
 K(x)=
 \frac{\sin(\pi x-1/\sqrt2)}{2\pi x-\sqrt2}
 +
 \frac{\sin(\pi x+1/\sqrt2)}{2\pi x+\sqrt2},
\]
with removable singularities interpreted by continuity.

\begin{proof}[Proof of Theorem~\ref{aiv:thm:IV}]
Apply Lemma~\ref{aiv:lem:poisson} to the infinite grid.  By
Lemma~\ref{aiv:lem:tail}, replacing the infinite grid by the finite central grid
changes the unnormalized sum by $O(L^{-2})$.  Since
$a_L\asymp1$, division by $a_LL^2$ makes this error $O(L^{-4})$.
The remaining infinite-grid expression, after normalization, is the overlap
in Lemma~\ref{aiv:lem:overlap}.  This gives $k(x-x')+o(1)$ uniformly on every
fixed compact set of normalized separations.
\end{proof}

\subsection{Audit notes and exact scope}

The proof above is independent of the zeta explicit formula, zero density,
and prime-side pair correlation.  It uses only:
\begin{enumerate}
\item Poisson summation on the critical lattice $2\pi\mathbb Z/L$;
\item compact support and regularity of the chosen window;
\item the fixed-width taper;
\item elementary integration by parts.
\end{enumerate}
Thus Input IV is proved independently from Poisson summation and elementary Fourier estimates for the exact support/taper convention used in the main manuscript.



\section{Certificate metadata}\label{app:cert}

\paragraph{Code and certificate availability.}
The complete Arb/FLINT verifier, the archived 256-bit certification run,
machine-readable parameters, execution logs, environment information,
SHA-256 hashes, and reproduction instructions are publicly available in the
reproducibility artifact~\cite{Hurtado26Artifact}.  The immutable snapshot
used for the present certificate is GitHub release
\texttt{v1.0.0-paper}, attached to commit \texttt{c57f53e}.  The release
contains the manuscript snapshot and the computational evidence used to
certify Proposition~\ref{prop:cert}.  Subsequent changes to the repository's
default branch are not part of the trust base of this certificate.

The frozen archival verifier is
\[
 \texttt{certify\_nonuniform\_3900\_v1\_1.py}.
\]
Its SHA-256 is
\begin{sloppypar}
\small\ttfamily
8eb7b4a17da8e114881264d3c2fc8daf262a0558d4248692b386c0fca2cc4301.
\end{sloppypar}
The 256-bit log SHA-256 is
\begin{sloppypar}
\small\ttfamily
4b0e62dcb5c4014504981f16e431144e3a01184b2aa2aa6b0bf650705d59db83.
\end{sloppypar}

\begin{center}
\begin{tabular}{@{}ll@{}}
\toprule
Field & Archived 256-bit run\\
\midrule
verified & \texttt{true}\\
target & $F_{6,\mathrm{nonuniform}}\ge39/10000$\\
grid & 4000\\
precision & 256 bits\\
nodes & 1,119,372\\
pruned & 560,334\\
splits & 559,038\\
maximum depth & 39\\
initial boxes & 1296\\
interval prunes & 364,112\\
tangent prunes & 190,953\\
pressure prunes & 5,269\\
elapsed time & 201.487482 s\\
\bottomrule
\end{tabular}
\end{center}

The pressure vector recorded by the verifier is
\[
 (2714,3733,3553,3553,3733,2714)/10^7.
\]
The kernel-table SHA-256 is
\begin{sloppypar}
\small\ttfamily
f75a3968f6daf7ea74867e430dc69fdca4276c4bb09f8cecc1b9846e8942c3f5.
\end{sloppypar}
The second-derivative-table SHA-256 is
\begin{sloppypar}
\small\ttfamily
f878d03d2eeca7fd34e3064a8dacb70d054fe3b3e9e3890c79a1f8297633cfb3.
\end{sloppypar}
The archived environment reports CPython 3.14.7 and
\texttt{python-flint==0.9.0}.

\section{Exact arithmetic}

From $m=262$ and $A_0=624/625$,
\[
 \frac{A_0}{m}=\frac{312}{81875},\qquad
 \frac{m-1}{500m}=\frac{261}{131000},\qquad
 1-\frac{312}{81875}=\frac{81563}{81875}.
\]
Thus
\[
 \frac{H_{\rm MT}-261/131000}{81563/81875}
 =
 \frac{655000H_{\rm MT}-1305}{652504}.
\]

\begin{thebibliography}{99}

\bibitem{Weil52}
A.~Weil,
\emph{Sur les ``formules explicites'' de la th\'eorie des nombres premiers},
Comm. S\'em. Math. Univ. Lund, Tome Suppl\'ementaire (1952), 252--265.

\bibitem{IK04}
H.~Iwaniec and E.~Kowalski,
\emph{Analytic Number Theory},
American Mathematical Society Colloquium Publications, Vol.~53, 2004.

\bibitem{MV74}
H.~L.~Montgomery and R.~C.~Vaughan,
\emph{Hilbert's inequality},
J. London Math. Soc. (2) \textbf{8} (1974), 73--82.

\bibitem{Titchmarsh86}
E.~C.~Titchmarsh,
\emph{The Theory of the Riemann Zeta-Function},
2nd ed., revised by D.~R.~Heath-Brown,
Oxford University Press, 1986.



\bibitem{Claude26}
Claude,
\emph{More than two thirds of the zeros of the Riemann zeta function lie on
the critical line},
preprint, August 10, 2026, Anthropic.

\bibitem{BGSTB24}
S.~A.~C.~Baluyot, D.~A.~Goldston, A.~I.~Suriajaya, and
C.~L.~Turnage-Butterbaugh,
\emph{An unconditional Montgomery theorem for pair correlation of zeros of
the Riemann zeta-function},
Acta Arith. \textbf{214} (2024), 357--376.

\bibitem{BGSTB25}
S.~A.~C.~Baluyot, D.~A.~Goldston, A.~I.~Suriajaya, and
C.~L.~Turnage-Butterbaugh,
\emph{Pair correlation of zeros of the Riemann zeta function I: proportions
of simple zeros and critical zeros},
arXiv:2501.14545 (2025).

\bibitem{GS25}
D.~A.~Goldston and A.~I.~Suriajaya,
\emph{Zeta zeros on the critical line},
arXiv:2511.20059v2 (2025).

\bibitem{GS26}
D.~A.~Goldston and A.~I.~Suriajaya,
\emph{Zeta zeros in a narrow vertical box},
arXiv:2603.28104 (2026).

\bibitem{Johansson17}
F.~Johansson,
\emph{Arb: efficient arbitrary-precision midpoint-radius interval arithmetic},
IEEE Trans. Comput. \textbf{66} (2017), no.~8, 1281--1292.

\bibitem{Mon73}
H.~L.~Montgomery,
\emph{The pair correlation of zeros of the zeta function},
in \emph{Analytic Number Theory (St. Louis, 1972)},
Proc. Sympos. Pure Math. \textbf{24}, Amer. Math. Soc., 1973, 181--193.

\bibitem{Mon75}
H.~L.~Montgomery,
\emph{Distribution of the zeros of the Riemann zeta function},
in \emph{Proceedings of the International Congress of Mathematicians
(Vancouver, 1974)}, Vol.~1, Canadian Mathematical Congress, 1975, 379--381.


\bibitem{Jo26}
S.~Jo (GitHub user \texttt{ainta}),
\emph{A 67.30085\% lower bound for simple zeros of the Riemann zeta
function},
version 0.1.0, 10 August 2026,
GitHub repository \texttt{ainta/zeta-simple-zeros}.
The repository reports use of GPT-5.6 Sol in generating the research draft.

\bibitem{Rademacher26}
L.~Rademacher,
\emph{AI Refines AI: a reproducible candidate improvement to the 67.25\%
zeta-zero bound},
version 0.2.0, 11 August 2026,
GitHub repository \texttt{learademacher/ai-refines-ai-zeta-bound}.
The repository records Jo's artifact as upstream provenance and reports use
of GPT-5.6 Sol in the refinement workflow.


\bibitem{Hurtado26Artifact}
M.~Hurtado,
\emph{Reproducibility artifact for ``A position-weighted refinement for
simple zeros of the Riemann zeta function''},
release \texttt{v1.0.0-paper}, 13 August 2026,
GitHub repository
\url{https://github.com/MichaelMobius/simple_zeros_of_the_riemann_zeta_function},
release snapshot
\url{https://github.com/MichaelMobius/simple_zeros_of_the_riemann_zeta_function/releases/tag/v1.0.0-paper},
commit \texttt{c57f53e}.

\end{thebibliography}

\end{document}
